\chapimg{fig/sh13.jpg}
\IfFileExists{extra/lead13.tex}{\chapterlead{\subhd{After this chapter you can} classify sprinkler systems; compute the nozzle discharge and the wetted radius; choose the spacing from the overlap rule; compute application rate, set time and the Christiansen uniformity; and size the lateral with the 20 per cent rule.}
}{}
\chapter{Sprinkler Irrigation Design}\label{chap:13}
\markboth{Sprinkler Irrigation Design}{Sprinkler Irrigation Design}
\IfFileExists{extra/bg13.tex}{\section{Introduction}
A sprinkler system imitates rain: water under pressure leaves through a nozzle, breaks into drops and falls on the crop. It can irrigate land that is not level, it can apply small depths, and it can be automated. Its drawbacks are energy (the pressure), wind drift and evaporation from the drops, and wetting of the leaves. This chapter gives the design from the nozzle to the lateral.

\section{Types of sprinkler systems}
Sprinkler systems are \emph{set} (portable or solid-set laterals moved between positions), \emph{continuous-move} (centre pivot, linear move, travelling gun) and \emph{permanent}. Solid-set systems need many sprinklers and pipes but little labour; pivots need little pipe but a smooth ground and an even application rate that rises towards the end.

\section{Nozzle discharge and wetted radius}
The discharge of a nozzle follows the orifice law $q=C_dA\sqrt{2gH}$. For $d=4.8$ mm and a pressure of 350 kPa with $C_d=0.95$, $q=1.64$ m$^3$\,h$^{-1}$. The wetted radius grows with the nozzle size and the pressure; an empirical form is $R=3.4\sqrt{d\,P}$ with $d$ in mm and $P$ in bar, which gives $R=13.9$ m here.

\section{Spacing and application rate}
With the sprinklers on a rectangular spacing $S_l\times S_m$ the application rate is $I=q/(S_lS_m)$; here $I=7.6$ mm\,h$^{-1}$, which must be below the infiltration rate of the soil, or runoff follows. The spacing is chosen at 50 to 65 per cent of the wetted diameter, less in wind.

\section{Uniformity}
Christiansen's coefficient is $CU=100\bigl(1-\sum|z_i-\bar z|/\sum z_i\bigr)$ from catch-can depths $z_i$; 85 per cent or more is the target. For the depths 11, 12, 9, 13, 10, 12, 8, 11 mm: mean 10.75, $\sum|z-\bar z|=10.50$, $CU=87.8$ per cent.

}{}
IfFileExists{extra/th13.tex}{\section{Types of sprinkler systems}
In \kw{sprinkler irrigation} the water is carried under pressure in pipes and sprayed through nozzles, so that it falls on the crop like rain. The systems are classified by their mobility.
\begin{exlist}
\item \kw{Hand-move lateral.} Aluminium pipes with couplers are moved by hand after each set. Low capital cost, high labour.
\item \kw{Solid-set.} The laterals and sprinklers stay in the field for the whole season (or are buried), and are controlled by valves; also used for frost protection and cooling.
\item \kw{Side-roll (wheel-line).} A lateral on large wheels is rolled to the next position; for low crops on rectangular fields.
\item \kw{Centre pivot.} A lateral on towers turns round a fixed centre, and irrigates a circle; one of the most widely used systems for large fields. The speed of the outer tower is greatest, so the nozzles are larger towards the end, with a discharge proportional to the radius.
\item \kw{Linear move.} The lateral moves in a straight line, and irrigates a rectangle.
\item \kw{Travelling gun (big gun).} One large sprinkler on a hose reel; flexible, but with high pressure and large drops.
\end{exlist}

\section{Components of a sprinkler system}
A pump and a power unit; the \emph{mainline} and the \emph{sub-main}, which can be buried; the \emph{lateral}, which carries the sprinklers; the \emph{riser}, which lifts the sprinkler above the crop; and the \emph{sprinkler head}, with one or two \emph{nozzles}. The head turns by the impact of a spring-loaded arm (\emph{impact} sprinkler) or by a gear drive. Small nozzles, at 200 to 300 kPa, give fine drops; large nozzles for big guns need 400 to 700 kPa.

\section{Advantages and disadvantages of sprinkler irrigation}
\textbf{Advantages:} it can be used on any slope without levelling; the application rate can be matched to the infiltration of a light soil; it gives an even distribution if the design is right; it can be automated; the depth is easy to control, and small depths can be given for germination, cooling and frost protection; it saves land (no channels).
\par\noindent \textbf{Disadvantages:} the energy for the pressure; a higher capital cost than surface irrigation; the wind distorts the pattern, and the drops can be blown away (drift loss); evaporation in the air; the wetting of leaves can promote disease and can scorch the leaves if the water is saline; and the large drops of a big gun can seal the surface of a bare soil.

\section{Drops, energy and the soil surface}
The drops that leave the nozzle have a diameter between 0.5 and 4 mm. A large drop falls faster and has more \kw{kinetic energy}; on a bare soil it breaks the aggregates, seals the surface and lowers the infiltration. A fine drop is blown by the wind and loses more by evaporation. The nozzle pressure, the nozzle size and the type of head control the spectrum of drops: a high pressure with a small nozzle gives fine drops, a low pressure with a large nozzle gives coarse ones. The aim is the largest drop that does not damage the soil and the crop.

\section{Wind and the pattern of a sprinkler}
Without wind a sprinkler gives a pattern that is roughly triangular in section, with the most water near the head, falling to nothing at the radius. Wind shifts the pattern downwind and shortens it upwind, and a wind above 3 m\,s$^{-1}$ requires closer spacing along the wind direction. The recommended spacing, as a fraction of the wetted diameter, is 65 per cent for a wind below 2 m\,s$^{-1}$, 60 per cent for 2 to 4 m\,s$^{-1}$ and 50 per cent above that. Overlapping patterns add up to an even depth, which the uniformity coefficient measures (Section~\ref{ch:cu}).

\section{Application rate and the infiltration of the soil}
The rate at which the sprinklers apply water must not be larger than the soil can take, or water will pond and run off. The soil takes water at a rate that falls to a final value, and the maximum allowable rate depends on the soil and the slope; typical limits are 20 to 25 mm\,h$^{-1}$ in a sand, 10 to 15 in a loam and 5 or less in a clay, reduced by 20 to 40 per cent on slopes. If the rate is too high, a smaller nozzle, a wider spacing or a different sprinkler must be chosen; the check is the last of the ten tests of this chapter.
}{}
\IfFileExists{extra/deep13.tex}{\section{Sprinkler irrigation: what the layer does}
The sprinkler layer simulates the distribution of water from a set of sprinklers over the field and the fate of that water in the soil. From a sprinkler type, a nozzle diameter, a pressure and a spacing, it computes the discharge and the wetted radius; it superposes the patterns of all the sprinklers around a point to get the depth caught at every point of a repeating cell; it removes the water lost to wind drift and evaporation; and it passes the rest through the intake of the soil to find the water stored, lost to deep percolation and lost as runoff. The set time is chosen so that the driest quarter of the field receives the net depth.

\section{Types of sprinklers}
\begin{table}[!htbp]\centering\footnotesize\setlength{\tabcolsep}{3pt}\caption{Sprinkler types of the application.}
\begin{tabular}{lrlrcccc}\toprule
\hd{Type} & \hd{$c_R$} & \hd{pattern} & \hd{drift $F$} & \hd{nozzle range (mm)} & \hd{pressure range (kPa)} & \hd{default $d$ / $P$} & \hd{default $S_s\times S_l$ (m)}\\\midrule
Impact sprinkler & 1.35 & 1.3 & 1.0 & 2.4--8.0 & 200--500 & 4.4 / 300 & 12 $\times$ 18\\
Rotating spray & 0.80 & ring & 1.6 & 1.5--5.0 & 70--250 & 3.2 / 170 & 6 $\times$ 6\\
Micro-sprinkler & 0.70 & ring2 & 0.7 & 0.8--2.5 & 100--250 & 1.3 / 180 & 3 $\times$ 3\\
Big gun & 0.95 & 1.0 & 1.3 & 10.0--25.0 & 350--700 & 16.0 / 500 & 30 $\times$ 36\\\bottomrule\end{tabular}\end{table}

\textbf{Impact sprinklers} are the classic rotating-arm heads, with one or two nozzles and a throw of 12 to 20 m; they give a triangular-like pattern and a good uniformity at moderate pressure. \textbf{Rotating sprays} throw less (5 to 8 m) at low pressure (70--250 kPa) and give a pattern that is heaviest in a ring; their small drops drift more (factor 1.6). \textbf{Micro-sprinklers} are small, low-pressure heads for trees, with a throw of 2 to 4 m and a flat pattern. \textbf{Big guns} are large rotating nozzles of 10 to 25 mm that throw 30 m or more at 350--700 kPa; they are used on travelling reels and are the most sensitive to wind. For each type the pattern is a function of the relative distance $x=r/R$ from the sprinkler:

\begin{equation}\hat p(x)=\begin{cases}1-x^{m}&\text{impact }(m=1.3),\ \text{big gun }(m=1)\\ (0.5+0.5x)(1-x^{6})&\text{rotating spray}\\ (0.55+0.45x)(1-x^{5})&\text{micro-sprinkler}\end{cases},\quad \hat p=0\ (x\ge1),\label{eq:p-pat}\end{equation}
\par\noindent where $\hat p(x)$ is the relative application rate at the relative distance $x=r/R$ (1 at the sprinkler), $r$ the distance from the sprinkler (m), $R$ the wetted radius (m) and $m$ the exponent of the pattern. The pattern is scaled so that its integral over the circle equals the discharge of the nozzle.

\section{Nozzle, pressure and radius}
\begin{equation}\boxed{q=0.95\,A_n\sqrt{2g\,h}\times6\times10^{4},\quad A_n=\pi\left(\frac{d}{2000}\right)^{2},\quad h=\frac{P}{9.81},}\label{eq:p-q}\end{equation}
\par\noindent where $q$ is the discharge of the nozzle (L\,min$^{-1}$), 0.95 the discharge coefficient, $A_n$ the area of the nozzle (m$^2$), $d$ its diameter (mm), $g=9.81$ m\,s$^{-2}$, $h$ the pressure head (m) and $P$ the operating pressure (kPa); $6\times10^{4}$ converts m$^3$\,s$^{-1}$ to L\,min$^{-1}$.

\begin{equation}\boxed{R=c_R\sqrt{d\,h}\;\Phi(\alpha),\qquad\Phi(\alpha)=\left[\frac{\max(0.05,\sin3.1\alpha^{\circ})}{0.9763}\right]^{0.45},}\label{eq:p-R}\end{equation}
\par\noindent where $R$ is the wetted radius (m), $c_R$ the radius coefficient of the type (Table above), $d$ the nozzle diameter (mm), $h$ the pressure head (m) and $\Phi(\alpha)$ a correction for the angle $\alpha$ of the trajectory above the horizontal (degrees); $\Phi=1$ at the default angle of 25$^\circ$, and the range peaks near $29^\circ$.

\exbox{Discharge and radius of a nozzle}{An impact sprinkler with a 4.4-mm nozzle at 300 kPa and a trajectory of 25$^\circ$: find the head, the discharge and the wetted radius.}\sol $h=300/9.81=30.58$ m \ans $A_n=\pi(4.4/2000)^2=1.521e-05$ m$^2$ \ans $q=0.95\times1.521e-05\times\sqrt{2\times9.81\times30.58}\times6\times10^4=21.23$ L\,min$^{-1}$ ($1.27$ m$^3$\,h$^{-1}$) \ans $\Phi(25^\circ)=1$; $R=1.35\sqrt{4.4\times30.58}=15.66$ m \ans Wetted diameter $31.3$ m; the spacing of 12$\times$18 m is 38 and 57 per cent of the diameter \finans

\begin{table}[!htbp]\centering\footnotesize\setlength{\tabcolsep}{3pt}\caption{Discharge and radius for three nozzle sizes of each type at its default pressure.}
\begin{tabular}{lrrrrrr}\toprule
\hd{Type} & \hd{$d$ (mm)} & \hd{$P$ (kPa)} & \hd{$h$ (m)} & \hd{$q$ (L/min)} & \hd{$R$ (m)} & \hd{$q$ (m$^3$/h)}\\\midrule
Impact sprinkler & 2.4 & 300 & 30.6 & 6.32 & 11.6 & 0.38\\
Impact sprinkler & 4.4 & 300 & 30.6 & 21.23 & 15.7 & 1.27\\
Impact sprinkler & 6.4 & 300 & 30.6 & 44.92 & 18.9 & 2.69\\
Rotating spray & 1.5 & 170 & 17.3 & 1.86 & 4.1 & 0.11\\
Rotating spray & 3.2 & 170 & 17.3 & 8.45 & 6.0 & 0.51\\
Rotating spray & 5 & 170 & 17.3 & 20.64 & 7.4 & 1.24\\
Micro-sprinkler & 0.8 & 180 & 18.3 & 0.54 & 2.7 & 0.03\\
Micro-sprinkler & 1.3 & 180 & 18.3 & 1.44 & 3.4 & 0.09\\
Micro-sprinkler & 2.5 & 180 & 18.3 & 5.31 & 4.7 & 0.32\\
Big gun & 10 & 500 & 51.0 & 141.57 & 21.4 & 8.49\\
Big gun & 16 & 500 & 51.0 & 362.41 & 27.1 & 21.74\\
Big gun & 25 & 500 & 51.0 & 884.80 & 33.9 & 53.09\\\bottomrule\end{tabular}\end{table}

\begin{table}[!htbp]\centering\footnotesize\setlength{\tabcolsep}{3pt}\caption{Effect of the pressure on the impact sprinkler of 4.4 mm.}
\begin{tabular}{rrrr}\toprule
\hd{$P$ (kPa)} & \hd{$h$ (m)} & \hd{$q$ (L/min)} & \hd{$R$ (m)}\\\midrule
200 & 20.4 & 17.33 & 12.8\\
250 & 25.5 & 19.38 & 14.3\\
300 & 30.6 & 21.23 & 15.7\\
350 & 35.7 & 22.93 & 16.9\\
400 & 40.8 & 24.51 & 18.1\\
500 & 51.0 & 27.41 & 20.2\\\bottomrule\end{tabular}\end{table}

\section{Overlap and uniformity}
The application rate at a point is the sum of the contributions of all sprinklers around it. The application evaluates it on a repeating cell of the layout, $S_s\times S_l$ for a rectangular layout and $S_s\times2S_l$ for a triangular layout (where every second lateral is shifted by $S_s/2$):

\begin{equation}\boxed{a(x,y)=\sum_{i,j}p\bigl(x-x_{ij},\,y-y_{ij}\bigr),\qquad \overline{a}=\frac{1}{n}\sum a,\qquad CU=1-\frac{\frac1n\sum|a-\bar a|}{\bar a},}\label{eq:p-cu}\end{equation}
\par\noindent where $a(x,y)$ is the gross application rate at the point $(x,y)$ (mm\,min$^{-1}$), $p$ the pattern of one sprinkler (Eq.~\ref{eq:p-pat} scaled to the discharge), $(x_{ij},y_{ij})$ the positions of the sprinklers (shifted by $S_s/2$ on alternate laterals in a triangular layout), $n$ the number of points of the cell (196 for a rectangular layout), $\bar a$ the mean rate and $CU$ the Christiansen uniformity coefficient (fraction); the mean rate over the cell equals the discharge divided by the area, $q/(S_sS_l)$.



The wind distorts the pattern: it shifts the centre downwind by $0.03uR$ and stretches the radius to $R(1+0.022u\cos\theta)$ downwind and shortens it upwind (limits 0.7 to 1.4). Here $u$ is the wind speed (m\,s$^{-1}$) and $\theta$ the angle from the wind direction.

\begin{table}[!htbp]\centering\footnotesize\setlength{\tabcolsep}{3pt}\caption{Spacing and layout for the impact sprinkler (4.4 mm, 300 kPa, wind 3 m/s) on the loam S06, $Z_n=12$ mm.}
\begin{tabular}{llrrrr}\toprule
\hd{Layout} & \hd{$S_s\times S_l$ (m)} & \hd{mean rate (mm/h)} & \hd{$CU$ (\%)} & \hd{$DU_{lq}$ (\%)} & \hd{set time (min)}\\\midrule
rectangular & 9 $\times$ 9 & 15.7 & 98.2 & 97.8 & 49\\
rectangular & 12 $\times$ 12 & 8.8 & 98.9 & 98.5 & 86\\
rectangular & 12 $\times$ 18 & 5.9 & 96.2 & 94.4 & 134\\
rectangular & 15 $\times$ 18 & 4.7 & 95.1 & 91.9 & 173\\
rectangular & 18 $\times$ 18 & 3.9 & 96.0 & 94.1 & 202\\
rectangular & 18 $\times$ 24 & 2.9 & 74.3 & 61.8 & 411\\
rectangular & 24 $\times$ 24 & 2.2 & 68.7 & 45.7 & 742\\
triangular & 9 $\times$ 9 & 15.7 & 98.5 & 98.2 & 48\\
triangular & 12 $\times$ 12 & 8.8 & 97.2 & 95.8 & 88\\
triangular & 12 $\times$ 18 & 5.9 & 97.9 & 97.2 & 131\\
triangular & 15 $\times$ 18 & 4.7 & 94.1 & 89.7 & 177\\
triangular & 18 $\times$ 18 & 3.9 & 93.5 & 89.8 & 212\\
triangular & 18 $\times$ 24 & 2.9 & 74.1 & 63.5 & 400\\
triangular & 24 $\times$ 24 & 2.2 & 69.5 & 56.1 & 604\\\bottomrule\end{tabular}\end{table}

\begin{table}[!htbp]\centering\footnotesize\setlength{\tabcolsep}{3pt}\caption{Effect of the wind on the impact sprinkler, 12 $\times$ 18 m, loam S06.}
\begin{tabular}{rrrrrr}\toprule
\hd{Wind $u$ (m/s)} & \hd{$CU$ (\%)} & \hd{$DU_{lq}$ (\%)} & \hd{drift loss (\%)} & \hd{set time (min)} & \hd{$E_a$ (\%)}\\\midrule
0 & 96.2 & 94.3 & 1.5 & 131 & 92.7\\
1 & 96.2 & 94.3 & 2.2 & 132 & 92.1\\
2 & 96.2 & 94.4 & 2.9 & 133 & 91.4\\
3 & 96.2 & 94.4 & 3.8 & 134 & 90.6\\
5 & 96.3 & 94.5 & 5.9 & 137 & 88.8\\
7 & 96.5 & 94.7 & 8.3 & 141 & 86.7\\
10 & 96.5 & 94.9 & 12.7 & 148 & 82.7\\\bottomrule\end{tabular}\end{table}

\section{Wind drift and evaporation}
\begin{equation}\delta=\min\Bigl\{0.45,\ \tfrac1{100}\bigl(1.2+0.5u+0.04u^{2}\bigr)F\Bigl(\tfrac{P}{250}\Bigr)^{0.7}\bigl[1+0.025(T_a-20)\bigr]\bigl[0.6+0.4\tfrac{\sin\alpha}{0.4226}\bigr]\Bigr\},\label{eq:p-wd}\end{equation}
\par\noindent where $\delta$ is the fraction of the discharge lost to drift and evaporation of the drops in the air (the percentage is clipped between 0.5 and 45), $u$ the wind speed at the sprinkler (m\,s$^{-1}$), $F$ the drift factor of the type (1, 1.6, 0.7, 1.3), $P$ the pressure (kPa), $T_a$ the air temperature ($^\circ$C) and $\alpha$ the trajectory angle (the last factor is 1 at 25$^\circ$).

\exbox{Drift loss}{Find the drift and evaporation loss of the impact sprinkler at 300 kPa, wind 3 m\,s$^{-1}$, 24$^\circ$C and a trajectory of 25$^\circ$.}\sol $1.2+0.5(3)+0.04(9)=3.06$ \ans $(P/250)^{0.7}=(1.2)^{0.7}=1.136$; $1+0.025(24-20)=1.10$; angle factor $0.6+0.4\sin25^\circ/0.4226=1.0$ \ans $\delta=3.06\times1\times1.136\times1.10\times1.0/100=0.0382$, i.e.\ 3.8 per cent \ans On a gross depth of 13.2 mm this is 0.51 mm lost before the water reaches the soil \finans

\section{The soil intake and the set time}
At a point that receives the rate $a$ (net of the drift) for the time $T$ the intake follows Kostiakov, $F=KT^{A}$ with $K$ in mm\,min$^{-A}$. Water arrives at the rate $a$ until the soil can no longer take it; ponding starts at

\begin{equation}t_p=\left(\frac{a}{K}\right)^{\frac1{A-1}},\qquad F=\begin{cases}aT&T\le t_p\\ KT^{A}&T>t_p\end{cases},\qquad RO=\max\{0,\ aT-F-S_{d,\rm eff}\},\quad S_{d,\rm eff}=\frac{S_d}{1+0.25S},\label{eq:p-ga}\end{equation}
\par\noindent where $t_p$ is the time of ponding (min), $a$ the application rate at the point (mm\,min$^{-1}$), $K$ (mm\,min$^{-A}$) and $A$ the Kostiakov parameters of the soil (from the Saxton--Rawls and Green--Ampt chain of Chapter~11), $T$ the time of the set (min), $F$ the infiltrated depth (mm), $RO$ the runoff (mm), $S_{d,\rm eff}$ the effective depression storage (mm), $S_d$ the surface storage (5, 3.5 and 2 mm for conventional, reduced and no tillage) and $S$ the ground slope (per cent).



The set time is the smallest $T$ for which the mean depth in the driest quarter of the points reaches the net depth $Z_n$ (a bisection over 1 to 2400 min). The performance indicators follow:

\begin{equation}\boxed{E_a=\frac{\overline{\min(I,Z_n)}}{\bar a T},\quad E_s=\frac{\overline{\min(I,Z_n)}}{Z_n},\quad DU_{lq}=\frac{\overline{I}_{lq}}{\bar I},}\label{eq:p-ind}\end{equation}
\par\noindent where $I$ is the depth infiltrated at a point (mm), $\bar a T$ the gross depth applied (mm), $\overline{\min(I,Z_n)}$ the mean stored depth over the points, $E_s$ the requirement efficiency and $DU_{lq}$ the distribution uniformity, the mean of the lowest quarter of the infiltrated depths $\overline I_{lq}$ divided by the mean $\bar I$; the deep percolation is $\overline{\max(I-Z_n,0)}/(\bar aT)$ and the runoff $\overline{RO}/(\bar aT)$.



\begin{table}[!htbp]\centering\footnotesize\setlength{\tabcolsep}{3pt}\caption{Impact sprinkler (4.4 mm, 300 kPa, 12 $\times$ 18 m, wind 3 m/s) on the ten soils, $Z_n=12$ mm.}
\begin{tabular}{llrrrrrrrr}\toprule
\hd{ID} & \hd{Soil} & \hd{$K_s$ (mm/h)} & \hd{$k$} & \hd{$A$} & \hd{set (min)} & \hd{$E_a$ (\%)} & \hd{$DP$ (\%)} & \hd{$RO$ (\%)} & \hd{$DU_{lq}$ (\%)}\\\midrule
S01 & Coarse loam & 13.1 & 0.0008 & 0.793 & 134 & 91 & 5.6 & 0.0 & 94\\
S02 & Coarse sand & 114.7 & 0.0056 & 0.827 & 134 & 91 & 5.6 & 0.0 & 94\\
S03 & Sand & 134.7 & 0.0055 & 0.853 & 134 & 91 & 5.6 & 0.0 & 94\\
S04 & Loamy sand & 86.4 & 0.0041 & 0.830 & 134 & 91 & 5.6 & 0.0 & 94\\
S05 & Sandy loam & 44.7 & 0.0031 & 0.778 & 134 & 91 & 5.6 & 0.0 & 94\\
S06 & Loam & 12.1 & 0.0017 & 0.693 & 134 & 91 & 5.6 & 0.0 & 94\\
S07 & Silt loam & 10.9 & 0.0020 & 0.671 & 134 & 91 & 5.6 & 0.0 & 94\\
S08 & Sandy clay loam & 5.3 & 0.0007 & 0.699 & 134 & 91 & 5.6 & 0.0 & 94\\
S09 & Clay loam & 3.2 & 0.0008 & 0.643 & 134 & 91 & 5.6 & 0.0 & 94\\
S10 & Silty clay loam & 3.8 & 0.0010 & 0.634 & 134 & 91 & 5.6 & 0.0 & 94\\\bottomrule\end{tabular}\end{table}

The table repeats the lesson of Chapter~6: the rate of the sprinkler, about 6 mm\,h$^{-1}$, is below the intake of most soils, so there is no runoff, and the main loss is the deep percolation from the over-irrigated overlap on the coarse soils; on the clays, where the intake is small, ponding and runoff begin sooner and the set time must be shortened.

\begin{table}[!htbp]\centering\footnotesize\setlength{\tabcolsep}{3pt}\caption{Effect of the pressure (12 $\times$ 18 m spacing is held, loam S06).}
\begin{tabular}{rrrrrr}\toprule
\hd{$P$ (kPa)} & \hd{$R$ (m)} & \hd{mean rate (mm/h)} & \hd{$CU$ (\%)} & \hd{set (min)} & \hd{$E_a$ (\%)}\\\midrule
200 & 12.8 & 4.8 & 82.2 & 199 & 75.0\\
250 & 14.3 & 5.4 & 92.4 & 157 & 84.9\\
300 & 15.7 & 5.9 & 96.2 & 134 & 90.6\\
350 & 16.9 & 6.4 & 96.6 & 126 & 89.2\\
400 & 18.1 & 6.8 & 95.7 & 121 & 87.3\\
450 & 19.2 & 7.2 & 95.5 & 113 & 87.6\\\bottomrule\end{tabular}\end{table}

\begin{table}[!htbp]\centering\footnotesize\setlength{\tabcolsep}{3pt}\caption{Effect of the nozzle size (300 kPa, 12 $\times$ 18 m, loam S06).}
\begin{tabular}{rrrrrrr}\toprule
\hd{$d$ (mm)} & \hd{$q$ (m$^3$/h)} & \hd{$R$ (m)} & \hd{rate (mm/h)} & \hd{$CU$ (\%)} & \hd{set (min)} & \hd{$E_a$ (\%)}\\\midrule
3.2 & 0.67 & 13.4 & 3.1 & 86.6 & 292 & 78.6\\
3.8 & 0.95 & 14.6 & 4.4 & 93.5 & 190 & 85.9\\
4.4 & 1.27 & 15.7 & 5.9 & 96.2 & 134 & 90.6\\
5.0 & 1.64 & 16.7 & 7.6 & 96.8 & 105 & 90.1\\
5.6 & 2.06 & 17.7 & 9.6 & 96.0 & 85 & 88.6\\
6.4 & 2.69 & 18.9 & 12.5 & 95.5 & 65 & 88.6\\\bottomrule\end{tabular}\end{table}

\exbox{A sprinkler set from nozzle to fate of the water}{Design the sprinkler set for the loam S06 with a net depth $Z_n=12$ mm: impact sprinkler 4.4 mm at 300 kPa, 12$\times$18 m, wind 3 m\,s$^{-1}$, air at 24$^\circ$C. Report the pattern, the set time and the fate of the water.}\sol \givenline{$Z_n=12$ mm}\ans Discharge $q=21.23$ L\,min$^{-1}$; radius $R=15.7$ m \ans Mean gross rate $\bar a=q/(S_sS_l)=21.23\times60/216=5.90$ mm\,h$^{-1}$ \ans Superposition on the $14\times14$ cell: $CU=96.2$ per cent \ans Drift loss $\delta=3.8$ per cent, so the net rate is $5.67$ mm\,h$^{-1}$ \ans The bisection finds $T=134$ min ($2.24$ h) for which the driest quarter catches 12 mm \ans Gross depth $13.2$ mm; $E_a=90.6$, $E_s=99.8$, deep percolation $5.6$, runoff $0.0$, drift $3.8$ per cent; $DU_{lq}=94.4$ per cent \finans

}{}
\banner{Sprinkler irrigation design}\label{ch:sprinkler}
\section{Nozzle discharge}
The pressure head of the nozzle is $h=P/9.81$ (m, $P$ in kPa) and the discharge follows the orifice law with coefficient 0.95,
\begin{equation}\boxed{q=0.95A_n\sqrt{2gh}\times6\times10^4\ \ (\mathrm{L\,min^{-1}}),\qquad A_n=\pi\left(\frac{d}{2000}\right)^2\ \mathrm{m^2}.}\end{equation}
\par\noindent where $q$ is the discharge of the nozzle (L\,min$^{-1}$), $A_n$ the area of the nozzle (m$^2$), $d$ its diameter (mm), $g$ the gravitational acceleration (m\,s$^{-2}$), $h$ the pressure head at the nozzle (m), 0.95 the discharge coefficient and $6\times10^4$ the factor from m$^3$\,s$^{-1}$ to L\,min$^{-1}$.

The factor $6\times10^4$ converts m$^3$ s$^{-1}$ to L min$^{-1}$. Because $q\propto d^2\sqrt P$, doubling the nozzle diameter quadruples the discharge, and doubling the pressure raises it by only 41\%.

\exbox{Discharge of one nozzle}{The default impact sprinkler has $d=4.4$ mm and $P=300$ kPa. Find $q$.}\sol \givenline{$d=4.4$ mm ; $P=300$ kPa}\ans $h=300/9.81=30.58$ m. $A_n=\pi(4.4/2000)^2=1.521\times10^{-5}$ m$^2$. $\sqrt{2gh}=\sqrt{2\times9.81\times30.58}=24.49$ m s$^{-1}$. So $q=0.95\times1.521\times10^{-5}\times24.49\times6\times10^4=21.23$ L min$^{-1}$ ($=1.274$ m$^3$ h$^{-1}$). \finans


\begin{table}[!htbp]\centering\small\caption{Discharge $q$ (L min$^{-1}$) of an impact nozzle, Eq.~(13.2).}\label{tab:qnoz}
\begin{tabular}{lrrrrr}\toprule
\hd{$d$ (mm)} & \hd{200 kPa} & \hd{250 kPa} & \hd{300 kPa} & \hd{400 kPa} & \hd{500 kPa}\\\midrule
2.4 & 5.2 & 5.8 & 6.3 & 7.3 & 8.2\\
3.2 & 9.2 & 10.3 & 11.2 & 13.0 & 14.5\\
4.0 & 14.3 & 16.0 & 17.5 & 20.3 & 22.7\\
4.4 & 17.3 & 19.4 & 21.2 & 24.5 & 27.4\\
5.0 & 22.4 & 25.0 & 27.4 & 31.7 & 35.4\\
6.0 & 32.2 & 36.0 & 39.5 & 45.6 & 51.0\\\bottomrule\end{tabular}\end{table}

\exbox{A lower discharge from the same nozzle}{Using Table~\ref{tab:qnoz}, a farmer wants to keep a 4.4 mm nozzle but cut the discharge by one quarter. What pressure is needed?}\sol  $q\propto\sqrt P$, so $q_2/q_1=\sqrt{P_2/P_1}=0.75$ gives $P_2/P_1=0.5625$, $P_2=0.5625\times300=169$ kPa. Check with the table: $q(4.4,169)$ is $15.92$ L min$^{-1}$, which is $75.0\%$ of 21.2. The price of the lower pressure is a smaller reach (next section) and a coarser drop. \finans


\section{Wetted radius}
The reach of a sprinkler is an empirical range law,
\begin{equation}\boxed{R=c_R\sqrt d\sqrt h\,f_\alpha(\alpha),\qquad f_\alpha=\left[\frac{\max\{0.05,\sin(3.1\alpha)\}}{0.9763}\right]^{0.45},}\end{equation}
with $d$ in mm, $h$ in m and the angle in degrees. The coefficient $c_R$ is 1.35 for the impact sprinkler (the Keller and Bliesner value), 0.80 for the rotating spray, 0.70 for the micro-sprinkler and 0.95 for the big gun; Table~\ref{tab:types} lists the defaults of the four types. The angle function equals 1 at $\alpha=25^\circ$ and peaks near $29^\circ$.

\begin{table}[!htbp]\centering\small\caption{The four sprinkler types at their defaults ($\alpha=25^\circ$): discharge and wetted radius.}\label{tab:types}
\begin{tabular}{lrrrrrrrr}\toprule
\hd{Type} & \hd{$d$} & \hd{$P$} & \hd{$S_s$} & \hd{$S_l$} & \hd{$c_R$} & \hd{$F$} & \hd{$q$ (L/min)} & \hd{$R$ (m)}\\\midrule
Impact & 4.4 & 300 & 12 & 18 & 1.35 & 1.0 & 21.2 & 15.7\\
Rotating spray & 3.2 & 170 & 6 & 6 & 0.80 & 1.6 & 8.5 & 6.0\\
Micro-sprinkler & 1.3 & 180 & 3 & 3 & 0.70 & 0.7 & 1.4 & 3.4\\
Big gun & 16.0 & 500 & 30 & 36 & 0.95 & 1.3 & 362.4 & 27.1\\\bottomrule\end{tabular}\end{table}

\exbox{Reach and trajectory angle}{Find the radius of the default impact sprinkler at $\alpha=25^\circ$ and at $\alpha=15^\circ$, and say which angle is the better for reach.}\sol \givenline{$\alpha=25^\circ$ ; $\alpha=15^\circ$}\ans $h=30.58$ m, $\sqrt d=\sqrt{4.4}=2.098$, $\sqrt h=5.530$. At $25^\circ$: $\sin77.5^\circ=0.9763$ so $f_\alpha=1$ and $R=1.35\times2.098\times5.530=15.66$ m. At $15^\circ$: $\sin46.5^\circ=0.7254$; $f_\alpha=(0.7254/0.9763)^{0.45}=0.8749$, so $R=13.70$ m. At $29^\circ$: $\sin89.9^\circ=0.99999$, $f_\alpha=1.0109$, $R=15.83$ m. The angle changes the reach by only a few per cent between $20^\circ$ and $30^\circ$, which is why the default $25^\circ$ is near the optimum, and a flat jet (low $\alpha$) shortens the reach but also lowers the drift (next sections). \finans


\begin{table}[!htbp]\centering\small\caption{Angle function, reach and drift of the impact sprinkler.}\label{tab:angle}
\begin{tabular}{lrrr}\toprule
\hd{$\alpha$ (deg)} & \hd{$f_\alpha$} & \hd{$R$ (m)} & \hd{$w_d$ (\%) at 3 m/s}\\\midrule
8 & 0.684 & 10.7 & 2.8\\
12 & 0.806 & 12.6 & 3.0\\
15 & 0.875 & 13.7 & 3.2\\
20 & 0.956 & 15.0 & 3.5\\
25 & 1.000 & 15.7 & 3.8\\
29 & 1.011 & 15.8 & 4.0\\
35 & 0.987 & 15.5 & 4.4\\
45 & 0.832 & 13.0 & 4.9\\\bottomrule\end{tabular}\end{table}

\begin{keybox}
At $\alpha=45^\circ$ the term $\sin(3.1\alpha)=\sin139.5^\circ=0.65$ has already fallen, and for $\alpha\gtrsim32^\circ$ the sine is declining: the angle function is a deliberate engineered shape, not a measured one. Table~\ref{tab:angle} shows the drift rising with the angle, since steeper jets stay longer in the air.
\end{keybox}
\section{The single-sprinkler pattern and wind}
With $x=r/R_e$ the shape of the deposit is
\begin{equation}\hat p(x)=\begin{cases}1-x^{1.3}&\text{impact}\\1-x&\text{big gun}\\(0.5+0.5x)(1-x^6)&\text{rotating spray}\\(0.55+0.45x)(1-x^5)&\text{micro-sprinkler}\end{cases}\qquad\hat p=0\ (x\ge1).\end{equation}
\par\noindent where $\hat p(x)$ is the relative application rate (fraction of the rate at the sprinkler) at the scaled distance $x=r/R$, where $r$ is the distance from the sprinkler and $R$ the wetted radius (m); the four forms are for impact sprinklers, big guns, rotating sprays and micro-sprinklers.

Wind stretches the radius downwind and moves the centre, $R_e(\theta)=R\,\mathrm{clip}\{1+0.022u\cos\theta;0.7,1.4\}$ and shift $=0.03uR$. The final pattern $p=s\hat p$ is scaled by a constant $s$ so that the volume deposited per minute equals the discharge, $\iint p\,dA=q$ (with $p$ in mm min$^{-1}$, since 1 mm over 1 m$^2$ is 1 L).

\exbox{Peak rate of a single sprinkler}{For no wind, find the exact peak application rate of the impact pattern, check it against the 45$\times$45 grid used by the layer, and give the rate at half the radius.}\sol  Volume of $\hat p=1-x^{m}$: $\iint\hat p\,dA=2\pi R^2\int_0^1(1-x^{m})x\,dx=2\pi R^2\left(\tfrac12-\tfrac1{m+2}\right)$. With $m=1.3$: $\tfrac12-\tfrac1{3.3}=0.19697$. $R=15.66$ m, so volume $=2\pi\times245.2\times0.19697=303.5$ m$^2$. The peak rate is $s=q/V=21.23/303.5=0.0700$ mm min$^{-1}$ $=4.20$ mm h$^{-1}$. The numerical grid gives $s=0.0699$ mm min$^{-1}$, differing by -0.0\% (the 45$\times$45 grid places the centre at a corner of four cells). At $x=0.5$: $p=0.0700\times(1-0.5^{1.3})=0.0415$ mm min$^{-1}$, i.e.\ 2.49 mm h$^{-1}$. \finans


\exbox{Wind distortion of the reach}{At wind speed $u=6.0$ m s$^{-1}$ find the downwind, upwind and crosswind radii of the impact pattern and the shift of its centre.}\sol \givenline{$u=6.0$ m s$^{-1}$ find}\ans Downwind ($\theta=0$): $R_e=15.66\times(1+0.022\times6.0)=15.66\times1.132=17.73$ m. Upwind ($\theta=180^\circ$): $R_e=15.66\times0.868=13.59$ m. Crosswind ($\theta=90^\circ$): $R_e=R=15.66$ m. The shift of the centre is $0.03\times6.0\times15.66=2.82$ m downwind. The elongation is +13.2\% downwind and $-$13.2\% upwind; the clip (0.7 to 1.4) would act only above $13.6$ m s$^{-1}$ upwind and $18.2$ m s$^{-1}$ downwind. \finans


\section{Overlap and the coefficient of uniformity}\label{ch:cu}
On a grid with spacings $S_s$ along the lateral and $S_l$ between laterals the rate at a point is the sum over neighbouring sprinklers, $a(x,y)=\sum_{ij}p(x-x_{ij},y-y_{ij})$, evaluated on a $14\times14$ grid of the repeating cell ($14\times24$ of the $S_s\times2S_l$ cell for the triangular layout, where each second lateral is shifted by $S_s/2$). The mean rate is $\mathrm{AR}=60\bar a$ mm h$^{-1}$ and the Christiansen coefficient is
\begin{equation}\boxed{\mathrm{CU}=1-\frac{\sum_k|a_k-\bar a|}{n\bar a}.}\end{equation}
\par\noindent where $CU$ is the Christiansen uniformity (0 to 1), $a_k$ the depth in can $k$ (mm), $\bar a$ the mean depth and $n$ the number of cans.

Because every drop of the discharge lands somewhere in the repeating cell, the mean rate does not depend on the pattern at all: $\mathrm{AR}=60q/(S_sS_l)$.

\exbox{Mean application rate}{Compute $\mathrm{AR}$ for the default impact layout ($12\times18$ m, $q=21.23$ L min$^{-1}$) and compare with the sum from the overlap grid.}\sol \givenline{$q=21.23$ L min$^{-1}$}\ans $\mathrm{AR}=60\times21.23/(12\times18)=5.897$ mm h$^{-1}$. The summed grid gives 5.898 mm h$^{-1}$ (wind 3 m s$^{-1}$), the same to within 0.1\%, which confirms the volume normalisation. The layer then compares this rate with the largest rate the soil allows (Section~\ref{sec:maxrate}). \finans


\begin{table}[!htbp]\centering\small\caption{Spacing along the lateral for the impact sprinkler ($S_l=18$ m, 3 m s$^{-1}$, $Z_n=30$ mm, loam).}\label{tab:ssweep}
\begin{tabular}{lrrrrr}\toprule
\hd{$S_s$ (m)} & \hd{$S_s/R$ (\%)} & \hd{AR (mm/h)} & \hd{CU (\%)} & \hd{$T$ (min)} & \hd{$E_a$ (\%)}\\\midrule
6 & 38 & 11.8 & 98.3 & 179 & 85.1\\
9 & 57 & 7.9 & 97.6 & 271 & 84.4\\
12 & 77 & 5.9 & 96.0 & 371 & 82.0\\
15 & 96 & 4.7 & 95.0 & 476 & 80.0\\
18 & 115 & 3.9 & 95.8 & 559 & 81.8\\
21 & 134 & 3.4 & 86.7 & 767 & 69.1\\
24 & 153 & 2.9 & 74.0 & 1{,}145 & 52.1\\\bottomrule\end{tabular}\end{table}

\exbox{Spacing and uniformity}{Using Table~\ref{tab:ssweep}, what is the widest spacing along the lateral that still gives CU $\ge85\%$, and does it satisfy the rule $S_s\le kR$ of the ten checks (Section~\ref{sec:checks})?}\sol  CU is 98.3\% at $S_s=6$ m, 96.0\% at 12 m, 95.8\% at 18 m and 74.0\% at 24 m. The last spacing at or above 85\% is 21 m. The wind factor for $u=3$ m s$^{-1}$ ($W=10.8$ km h$^{-1}$) is $k=0.9$, so the rule allows $S_s\le0.9\times15.7=14.1$ m. A pattern can therefore keep CU above 85\% at a spacing the rule rejects, or the reverse; the layer reports both and lets the person decide. \finans


\begin{table}[!htbp]\centering\small\caption{Square against triangular layout, impact sprinkler at $12\times18$ m, 3 m s$^{-1}$.}\label{tab:tri}
\begin{tabular}{lrrrrr}\toprule
\hd{Layout} & \hd{AR} & \hd{CU (\%)} & \hd{DU (\%)} & \hd{$T$ (min)} & \hd{$E_a$ (\%)}\\\midrule
Square & 5.90 & 96.0 & 93.4 & 371 & 82.0\\
Triangular & 5.90 & 97.7 & 96.6 & 360 & 84.6\\\bottomrule\end{tabular}\end{table}

\begin{table}[!htbp]\centering\small\caption{Effect of wind on the overlapped pattern of the default impact layout.}\label{tab:wind}
\begin{tabular}{lrrr}\toprule
\hd{$u$ (m/s)} & \hd{AR} & \hd{CU (\%)} & \hd{$w_d$ (\%)}\\\midrule
0 & 5.90 & 96.0 & 1.5\\
3 & 5.90 & 96.0 & 3.8\\
6 & 5.90 & 96.3 & 7.0\\
9 & 5.90 & 96.5 & 11.2\\
12 & 5.95 & 96.6 & 16.2\\\bottomrule\end{tabular}\end{table}

\begin{keybox}
In Table~\ref{tab:wind} the computed CU does \emph{not} fall as the wind rises: in this model wind only stretches and shifts every pattern in the same direction, so the superposition stays nearly uniform on a repeating cell, while the drift loss $w_d$ and the run time grow. Field experience says the opposite (wind spoils uniformity). The layer's wind rule therefore enters through the spacing limit $S_s\le kR$ and the drift loss, not through CU, and this is a known limitation of the layer.
\end{keybox}
\section{Drift and evaporation loss}
\begin{equation}\boxed{w_d=\frac1{100}\mathrm{clip}\Big\{(1.2+0.5u+0.04u^2)\,F\left(\frac P{250}\right)^{0.7}\![1+0.025(T_a-20)]\Big[0.6+0.4\frac{\sin\alpha}{0.4226}\Big];0.5,45\Big\}}\end{equation}
where $F$ is the type factor (impact 1.0, spray 1.6, micro 0.7, gun 1.3), $P$ in kPa, $u$ in m s$^{-1}$ and $T_a=24^\circ$C$+\Delta T$. The form is plausible (loss rises with wind, pressure and temperature, and falls with a flatter jet), but the coefficients are engineered, not fitted.

\exbox{Drift loss}{Find $w_d$ for the default impact sprinkler at $u=3$ m s$^{-1}$, $T_a=24^\circ$C.}\sol \givenline{$u=3$ m s$^{-1}$ ; $T_a=24^\circ$ C}\ans Wind term $1.2+0.5\times3+0.04\times9=3.06$. Pressure term $(300/250)^{0.7}=1.1361$. Temperature term $1+0.025\times4=1.1$. Angle term $0.6+0.4\sin25^\circ/0.4226=0.6+0.4\times0.4226/0.4226=1.000$ (since $\sin25^\circ=0.4226$ the bracket is exactly 1). Then $3.06\times1.0\times1.1361\times1.1\times1=3.824$ per cent, within the clip, so $w_d=0.0382$ (3.82\%). Compare the same sprinkler at 9 m s$^{-1}$: $1.2+4.5+3.24=8.94$, $w_d=11.2\%$. \finans


\begin{table}[!htbp]\centering\small\caption{Drift loss $w_d$ (\%) of the four types at their default pressures, $\alpha=25^\circ$, $T_a=24^\circ$C.}\label{tab:wd}
\begin{tabular}{lrrrr}\toprule
\hd{Type} & \hd{0 m/s} & \hd{3 m/s} & \hd{6 m/s} & \hd{9 m/s}\\\midrule
Impact & 1.5 & 3.8 & 7.0 & 11.2\\
Rotating spray & 1.6 & 4.1 & 7.6 & 12.0\\
Micro-sprinkler & 0.7 & 1.9 & 3.5 & 5.5\\
Big gun & 2.8 & 7.1 & 13.1 & 20.8\\\bottomrule\end{tabular}\end{table}

\section{Intake, ponding and runoff}\label{sec:pond}
The soil takes water by the Kostiakov law $F=Kt^A$ with $K=1000k$ mm min$^{-A}$ (Chapter~\ref{ch:surf}). At a point receiving $a_e=a(1-w_d)$ mm min$^{-1}$ the surface ponds when the rate exceeds the intake rate, $a_e=Kt^{A-1}$, which gives
\begin{equation}t_p=(a_e/K)^{1/(A-1)},\qquad F=\begin{cases}\mathrm{sup}=a_eT&T\le t_p\\KT^A&T>t_p\end{cases},\quad \mathrm{RO}=\max\{0,\,\mathrm{sup}-F-S_{d,\rm eff}\},\ I=\mathrm{sup}-\mathrm{RO}.\end{equation}
\par\noindent where $t_p$ is the time of ponding (h), $a_e$ the application rate (mm\,h$^{-1}$), $K$ and $A$ the parameters of the Kostiakov infiltration law $F=KT^A$ ($0<A<1$), $T$ the time (h), ${\rm sup}$ the supplied depth (mm), $F$ the infiltrated depth (mm), ${\rm RO}$ the runoff (mm), $S_{d,{\rm eff}}$ the effective surface storage (mm) and $I$ the infiltrated depth.

Surface storage is reduced by slope, $S_{d,\rm eff}=\max(0,S_d)/(1+0.25S)$ ($S$ in \%). The canopy and mulch keep $I_c=\min\{I,I_{\rm int}\}$, $I_{\rm int}=0.6\,\mathrm{CAP}[\text{mulch}]+1.2\,\mathrm{LAI}$, which evaporates and never reaches the soil. The depth that reaches the soil is $Z_k=I-I_c$.

\exbox{Ponding time on the loam}{Conventional tillage gives $S_d=5$ mm on a 1\% slope. A loam under maize with LAI = 2.5 and no mulch is irrigated at the mean rate of the default impact layout. Find $S_{d,\rm eff}$, $I_{\rm int}$, the intake coefficients and the ponding time.}\sol \givenline{$S_d=5$ mm}\ans $S_{d,\rm eff}=5/(1+0.25\times1)=4.00$ mm. $I_{\rm int}=1.2\times2.5=3.0$ mm (no mulch term). The loam's Kostiakov law (Chapter~\ref{ch:surf}) has $k=0.001675$ m min$^{-A}$ and $A=0.694$, so $K=1.675$ mm min$^{-A}$. The applied rate is $a_e=0.0983\,(1-0.0382)=0.0945$ mm min$^{-1}$ ($5.67$ mm h$^{-1}$ after drift). Then $t_p=(0.0945/1.675)^{1/(0.694-1)}=0.0564^{-3.271}=12{,}118$ min. A set is about 371 min long, so $T<t_p$ and the whole supply infiltrates: RO = 0. The loam takes this sprinkler with a wide margin. \finans


\begin{table}[!htbp]\centering\small\caption{Default impact and rotating-spray layouts on the ten soils ($Z_n=30$ mm, 1\% slope, $u=3$ m s$^{-1}$, $S_d=4$ mm effective, LAI 2.5). Max is the allowed rate of Table~\ref{tab:maxrate}.}\label{tab:sprsoils}
\begin{tabular}{llrrrrrrrr}\toprule
\hd{Soil} & \hd{Sprinkler} & \hd{$K$} & \hd{$A$} & \hd{AR} & \hd{Max} & \hd{$t_p$} & \hd{$T$} & \hd{RO\%} & \hd{$E_a$\%}\\\midrule
Coarse loam & Impact & 0.76 & 0.805 & 5.9 & 25 & 43{,}886 & 371 & 0.0 & 82.0\\
Coarse loam & Rotating spray & 0.76 & 0.805 & 14.1 & 25 & 517 & 209 & 3.1 & 60.4\\
Coarse sand & Impact & 5.32 & 0.837 & 5.9 & 50 & $>10^5$ & 371 & 0.0 & 82.0\\
Coarse sand & Rotating spray & 5.32 & 0.837 & 14.1 & 50 & $>10^5$ & 209 & 0.0 & 60.4\\
Sand & Impact & 5.28 & 0.862 & 5.9 & 50 & $>10^5$ & 371 & 0.0 & 82.0\\
Sand & Rotating spray & 5.28 & 0.862 & 14.1 & 50 & $>10^5$ & 209 & 0.0 & 60.4\\
Loamy sand & Impact & 3.91 & 0.840 & 5.9 & 50 & $>10^5$ & 371 & 0.0 & 82.0\\
Loamy sand & Rotating spray & 3.91 & 0.840 & 14.1 & 50 & $>10^5$ & 209 & 0.0 & 60.4\\
Sandy loam & Impact & 2.91 & 0.789 & 5.9 & 25 & $>10^5$ & 371 & 0.0 & 82.0\\
Sandy loam & Rotating spray & 2.91 & 0.789 & 14.1 & 25 & $>10^5$ & 209 & 0.0 & 60.4\\
Loam & Impact & 1.67 & 0.694 & 5.9 & 13 & 12{,}118 & 371 & 0.0 & 82.0\\
Loam & Rotating spray & 1.67 & 0.694 & 14.1 & 13 & 709 & 209 & 0.1 & 60.4\\
Silt loam & Impact & 1.92 & 0.669 & 5.9 & 13 & 8{,}878 & 371 & 0.0 & 82.0\\
Silt loam & Rotating spray & 1.92 & 0.669 & 14.1 & 13 & 647 & 209 & 0.1 & 60.4\\
Sandy clay loam & Impact & 0.69 & 0.701 & 5.9 & 8 & 776 & 371 & 0.0 & 82.0\\
Sandy clay loam & Rotating spray & 0.69 & 0.701 & 14.1 & 8 & 43 & 223 & 30.6 & 56.6\\
Clay loam & Impact & 0.79 & 0.636 & 5.9 & 4 & 344 & 371 & 0.0 & 82.0\\
Clay loam & Rotating spray & 0.79 & 0.636 & 14.1 & 4 & 32 & 291 & 47.3 & 43.8\\
Silty clay loam & Impact & 1.03 & 0.627 & 5.9 & 4 & 598 & 371 & 0.0 & 82.0\\
Silty clay loam & Rotating spray & 1.03 & 0.627 & 14.1 & 4 & 58 & 224 & 31.0 & 56.4\\\bottomrule\end{tabular}\end{table}

\begin{keybox}
Table~\ref{tab:sprsoils} shows where the ponding rule bites: runoff appears only on the soils whose intake coefficient is small enough that $t_p$ falls below the run time. In 4 of the 20 soil--sprinkler pairs the run produces runoff above 0.5\% of the gross depth. The allowed-rate column is a second, independent screen from the sprinkler catalogues; the two screens do not always agree, because Table~\ref{tab:maxrate} is a tabulated rule of thumb while $t_p$ uses the soil's computed intake.
\end{keybox}
\exbox{A partition that closes on the worst soil}{Take the soil with the most runoff in Table~\ref{tab:sprsoils} (Clay loam) under the rotating spray. Show how the water partition closes.}\sol  AR = 14.1 mm h$^{-1}$ against the table limit 4.0 mm h$^{-1}$; $w_d=4.1\%$; run time $T=291$ min (automatic); $t_p=32$ min. The partition of the gross depth 68.4 mm: $E_a=43.8\%$, deep percolation $0.4\%$, runoff $47.3\%$, drift plus interception $8.5\%$. Sum: $43.8+0.4+47.3+8.5=100.0\%$. The partition always closes because every millimetre discharged is either lost to drift, retained by the canopy, run off, stored, or drained. \finans


\section{Efficiencies and run time}
With $Z_k$ the depth reaching the soil at cell $k$ and $\mathrm{gross}=\bar aT$ the mean depth discharged,
\begin{equation}E_a=\frac{\overline{\min(Z_k,Z_n)}}{\mathrm{gross}},\ E_s=\frac{\overline{\min(Z_k,Z_n)}}{Z_n},\ \mathrm{DU}=\frac{Z_{25}}{\bar Z},\ \mathrm{DP}=\frac{\overline{\max(Z_k-Z_n,0)}}{\mathrm{gross}},\ \mathrm{RO}=\frac{\overline{\mathrm{RO}_k}}{\mathrm{gross}},\ \mathrm{WD}=w_d+\frac{\overline{I_c}}{\mathrm{gross}},\end{equation}
where $Z_{25}$ is the mean of the lowest quarter of $Z_k$. The automatic run time is the smallest $T\in[1,2400]$ min with $Z_{25}(T)\ge Z_n$, found by 18 bisection steps.

\exbox{The water partition of the default layout}{For the default layout on the loam ($Z_n=30$ mm) the layer finds $T=371$ min. Verify the efficiencies by hand from the summary of the cells.}\sol \givenline{$Z_n=30$ mm ; $T=371$ min}\ans Gross depth $=\bar aT=0.0983\times371=36.51$ mm ($\bar a=5.90$ mm h$^{-1}$). After drift ($w_d=3.8\%$) the supply is $35.12$ mm and the canopy retains $I_c=3.0$ mm, leaving a mean $\bar Z=32.12$ mm at the soil. The lowest-quarter mean is $Z_{25}=30.00$ mm, equal to $Z_n$ as required by Eq.~(13.12), so the run was solved correctly. Then $\mathrm{DU}=30.00/32.12=93.4\%$. Because every cell receives at least about $Z_n$, $E_s=99.8\%$ and $E_a=\overline{\min}/\mathrm{gross}=82.0\%$. DP $=6.0\%$, WD $=3.8+3.0/36.51=12.0\%$ and RO $=0$. Sum: $82.0+6.0+0+12.0=100.0\%$. \finans


\begin{table}[!htbp]\centering\small\caption{The four sprinklers on the loam at their default settings ($Z_n=30$ mm, 3 m s$^{-1}$).}\label{tab:fourtypes}
\begin{tabular}{lrrrrrrrrr}\toprule
\hd{Type} & \hd{$q$} & \hd{$R$} & \hd{AR} & \hd{CU} & \hd{$w_d$} & \hd{$T$} & \hd{$E_a$} & \hd{DP} & \hd{WD}\\\midrule
Impact & 21.2 & 15.7 & 5.9 & 96.0 & 3.8 & 371 & 82.0 & 6.0 & 12.0\\
Rotating spray & 8.5 & 6.0 & 14.1 & 77.7 & 4.1 & 209 & 60.4 & 29.3 & 10.2\\
Micro-sprinkler & 1.4 & 3.4 & 9.6 & 87.7 & 1.9 & 250 & 75.1 & 15.5 & 9.4\\
Big gun & 362.4 & 27.1 & 20.1 & 85.2 & 7.1 & 133 & 66.8 & 19.3 & 13.8\\\bottomrule\end{tabular}\end{table}

\exbox{Comparing the four types}{Using Table~\ref{tab:fourtypes}, which type gives the best $E_a$ on the loam, and what limits the worst?}\sol  The best is the Impact ($E_a=82.0\%$, CU $96.0\%$); the worst is the Rotating spray ($E_a=60.4\%$, CU $77.7\%$). For the worst, the pattern of a rotating spray is a ring that peaks away from the centre ($\hat p=(0.5+0.5x)(1-x^6)$ with a reach $R$ of only about 6 m against a 6 m spacing in both directions), so the overlap is incomplete; its low CU turns into deep percolation under the wet spots (DP $=29.3\%$) because the driest quarter must still receive $Z_n$. Uniformity, not the drift, controls $E_a$ here. \finans


\section{Maximum allowable application rate}\label{sec:maxrate}
The allowed rate (mm h$^{-1}$) is read from Table~\ref{tab:maxrate} by soil group and slope class; the ten soils are assigned to rows (coarse sand, sand, loamy sand $\to$ 1; sandy loam, coarse loam $\to$ 3; loam, silt loam $\to$ 5; sandy clay loam $\to$ 6; clay loam, silty clay loam $\to$ 7), and a compact layer moves an odd row to the next one.

\begin{table}[!htbp]\centering\small\caption{Maximum application rate (mm h$^{-1}$) by soil and slope (Keller and Bliesner).}\label{tab:maxrate}
\begin{tabular}{lrrrr}\toprule
\hd{Soil group} & \hd{0--5\%} & \hd{5--8\%} & \hd{8--12\%} & \hd{12--16\%}\\\midrule
Deep coarse sandy soils & 50 & 38 & 25 & 13\\
Coarse sandy, compact layer below & 38 & 25 & 19 & 10\\
Deep sandy loam & 25 & 20 & 15 & 10\\
Sandy loam, compact layer below & 19 & 13 & 10 & 8\\
Deep silt loam & 13 & 10 & 8 & 5\\
Silt loam, compact layer below & 8 & 6 & 4 & 2.5\\
Clay or clay loam, strong structure & 4 & 2.5 & 2 & 1.5\\\bottomrule\end{tabular}\end{table}

\exbox{Looking up the allowed rate}{Read the allowed rate for five cases.}\sol \begin{itemize}
\item Loam, slope 1\: row 5 ``Deep silt loam'', column 1: 13.0 mm h$^{-1}$.
\item Loam, slope 6\: row 5 ``Deep silt loam'', column 2: 10.0 mm h$^{-1}$.
\item Loam, slope 6\, compact layer: row 6 ``Silt loam, compact layer below'', column 2: 6.0 mm h$^{-1}$.
\item Clay loam, slope 2\: row 7 ``Clay or clay loam, strong structure'', column 1: 4.0 mm h$^{-1}$.
\item Coarse sand, slope 10\: row 1 ``Deep coarse sandy soils'', column 3: 25.0 mm h$^{-1}$.
\end{itemize}
The default impact layout has AR $=5.9$ mm h$^{-1}$ (and the loam at 1\% allows 13), so it passes. The same loam on a 6\% slope with a compact layer drops to 6 mm h$^{-1}$: the default 5.9 mm h$^{-1}$ still passes, but the rotating spray (AR $=14.1$) fails. \finans


\section{The ten textbook checks}\label{sec:checks}
With $h_s=P/9.81$, $q_h=0.06q$ (m$^3$ h$^{-1}$), $W=3.6u$ (km h$^{-1}$) and wind factor $k=1,0.9,0.8,0.7$ for $W\le5.5$, $\le13$, $\le22$, $>22$:
\begin{equation}\boxed{S_s\le kR,\quad S_l\le1.5R,\quad L_y=\frac{10^4A_f}{L_x},\quad N=\lfloor L_x/S_s\rfloor,\quad N_m=\lfloor L_y/S_l\rfloor,}\end{equation}
\par\noindent where $S_s$ and $S_l$ are the spacings of the sprinklers and the laterals (m), $R$ the wetted radius (m), $k$ the spacing factor (0.5 to 0.65, less in wind), $L_x$ and $L_y$ the lengths of the field sides (m), $A_f$ the area of the field (ha), $N$ the number of sprinklers per lateral and $N_m$ the number of lateral positions.

\begin{equation}t_{\rm set}=\frac{Z_{\rm gross}}{\mathrm{AR}},\quad n_{\rm sets/day}=\left\lfloor\frac{h_{\rm op}}{t_{\rm set}}\right\rfloor,\quad T_{\rm int}=\max\left(1,\left\lfloor\frac{Z_n}{ET_{\rm pk}}\right\rfloor\right),\quad N_d(n)=\left\lceil\frac{N_m}{n\,n_{\rm sets/day}}\right\rceil,\quad Q_{\rm in}(n)=nNq_h.\end{equation}
\par\noindent where $t_{\rm set}$ is the time of one set (h), $Z_{\rm gross}$ the gross depth (mm), ${\rm AR}$ the application rate (mm\,h$^{-1}$), $n_{{\rm sets/day}}$ the number of sets per day, $h_{\rm op}$ the hours of operation per day, $T_{\rm int}$ the irrigation interval (d), $Z_n$ the net depth (mm), $ET_{\rm pk}$ the peak crop water use (mm\,d$^{-1}$), $N_d(n)$ the days to cover the field with $n$ laterals, $N_m$ the number of lateral positions, $Q_{\rm in}(n)$ the inflow (L\,h$^{-1}$) and $q_h$ the discharge of one sprinkler.

The program picks the smallest $n$ ($1\ldots4$ laterals at once) with $N_d\le T_{\rm int}$ and $Q_{\rm in}\le3.6Q_s$ (m$^3$ h$^{-1}$, $Q_s$ in L s$^{-1}$).

\exbox{How many laterals run at once}{Farm: $A_f=5$ ha, $Q_s=40$ L s$^{-1}$, $L_L=150$ m, $h_{\rm op}=16$ h d$^{-1}$, $ET_{\rm pk}=6.4$ mm d$^{-1}$, $Z_n=30$ mm; default impact layout, $u=3$ m s$^{-1}$. Find the layout and the number of laterals to run at once.}\sol \givenline{$A_f=5$ ha ; $Q_s=40$ L s$^{-1}$ ; $L_L=150$ m ; $h_{\rm op}=16$ h d$^{-1}$ ; $ET_{\rm pk}=6.4$ mm d$^{-1}$ ; $Z_n=30$ mm; default impact layout ; $u=3$ m s$^{-1}$}\ans $h_s=30.58$ m, $q_h=0.06\times21.23=1.274$ m$^3$ h$^{-1}$, $W=10.8$ km h$^{-1}$ so $k=0.9$. Spacing: $S_s=12\le0.9\times15.66=14.1$ (OK) and $S_l=18\le1.5\times15.66=23.5$ (OK). $L_y=10^4\times5/150=333.3$ m; $N=\lfloor150/12\rfloor=12$ sprinklers per lateral; $N_m=\lfloor333.3/18\rfloor=18$ lateral positions. $t_{\rm set}=36.51/5.90=6.19$ h, so $n_{\rm sets/day}=\lfloor16/6.19\rfloor=2$; $T_{\rm int}=\lfloor30/6.4\rfloor=4$ d. Table~\ref{tab:nlat} gives $N_d$ and $Q_{\rm in}$; the supply limit is $3.6\times40=144$ m$^3$ h$^{-1}$. The smallest $n$ meeting both is $n=3$. \finans


\begin{table}[!htbp]\centering\small\caption{Laterals running together for the farm of the worked example.}\label{tab:nlat}
\begin{tabular}{lrrrr}\toprule
\hd{$n$} & \hd{$N_d$ (days)} & \hd{$Q_{\rm in}$ (m$^3$/h)} & \hd{$N_d\le T_{\rm int}$} & \hd{$Q_{\rm in}\le144$}\\\midrule
1 & 9 & 15.3 & no & yes\\
2 & 5 & 30.6 & no & yes\\
3 & 3 & 45.9 & yes & yes\\
4 & 3 & 61.1 & yes & yes\\\bottomrule\end{tabular}\end{table}

\section{Pipes, head and pump (Blasius form)}
For smooth plastic pipe the friction gradient is $J=0.472Q^{1.75}/D^{4.75}$ (m m$^{-1}$, $Q$ in L h$^{-1}$, $D$ in mm). With the Christiansen factor $F=1/2.75+1/(2N)$ for $N$ equal outlets and a 10\% allowance for fittings, the lateral takes the smallest internal diameter in $\{46.4,58,63,65.8,75,84.9,100\}$ mm for which
\begin{equation}\boxed{h_{f,\rm lat}=1.1FJ(Q_L,D)(L_x-S_s/2)\le0.2h_s,\qquad Q_L=1000Nq_h.}\end{equation}
\par\noindent where $h_{f,{\rm lat}}$ is the head loss of the lateral (m), $F$ the Christiansen outlet factor, $J(Q_L,D)$ the Hazen--Williams friction gradient (m\,m$^{-1}$) at flow $Q_L$ (L\,h$^{-1}$) and diameter $D$ (mm), $L_x$ the lateral length (m), $S_s$ the sprinkler spacing (m), $h_s$ the operating head of the sprinkler (m), 1.1 an allowance for fittings and 0.2 the 20 per cent rule.

The mainline has $Q_m=nQ_L$ and length $L_y-S_l/2$ and takes the smallest diameter not below the lateral's with $J<0.05$ m m$^{-1}$. The head and the pump are
\begin{equation}H_{\rm man}=z_{\rm lift}+JL_{\rm conv}+JL_{\rm main}+h_s+0.75h_{f,\rm lat},\qquad P_{\rm CV}=\frac{1000(Q_m/3.6\times10^6)H_{\rm man}}{75\eta},\ \eta=0.75,\ P_{\rm kW}=0.7355P_{\rm CV}.\end{equation}
\par\noindent where $H_{\rm man}$ is the head at the pump manifold (m), $z_{\rm lift}$ the lift (m), $J$ the friction gradient, $L_{\rm conv}$ and $L_{\rm main}$ the lengths of the conveyance and the mainline (m), $h_s$ the sprinkler head, $h_{f,{\rm lat}}$ the lateral loss, $P_{\rm CV}$ the power in metric horsepower, $Q_m$ the mainline flow (L\,h$^{-1}$), $\eta$ the pump and motor efficiency and $P_{\rm kW}$ the power in kW.


\begin{table}[!htbp]\centering\small\caption{Lateral loss for each standard diameter, $Q_L=15{,}285$ L h$^{-1}$, $N=12$ outlets.}\label{tab:latD}
\begin{tabular}{lrrr}\toprule
\hd{$D$ (mm)} & \hd{$J$ (m/m)} & \hd{$h_f$ (m)} & \hd{$h_f/h_s$ (\%)}\\\midrule
46.4 & 0.12036 & 7.73 & 25.3\\
58.0 & 0.04170 & 2.68 & 8.8\\
63.0 & 0.02816 & 1.81 & 5.9\\
65.8 & 0.02290 & 1.47 & 4.8\\
75.0 & 0.01230 & 0.79 & 2.6\\
84.9 & 0.00683 & 0.44 & 1.4\\
100.0 & 0.00314 & 0.20 & 0.7\\\bottomrule\end{tabular}\end{table}

\exbox{Lateral, mainline and pump}{Size the lateral, mainline and pump of the farm above.}\sol  $Q_L=1000\times12\times1.274=15{,}285$ L h$^{-1}$. The Christiansen factor is $F=1/2.75+1/(2\times12)=0.4053$. The allowed loss is $0.2h_s=0.2\times30.58=6.12$ m. Table~\ref{tab:latD} shows the first diameter with $h_f\le6.12$ m is $D=58.0$ mm ($h_f=2.68$ m, 8.8\% of $h_s$). Mainline: $Q_m=3\times15{,}285=45{,}856$ L h$^{-1}$; the first diameter not below 58.0 mm with $J<0.05$ is $D=84.9$ mm ($J=0.04667$ m m$^{-1}$). Head: $H_{\rm man}=12+0.04667\times120+0.04667\times324.3+30.58+0.75\times2.68=65.33$ m. Pump: $P_{\rm CV}=1000\times(45{,}856/3.6\times10^6)\times65.33/(75\times0.75)=14.79$ CV $=10.88$ kW. \finans


\section{The design table and sensitivity}
The optimiser maximises the score $\min(E_a,1)-2.5\max(0,0.9-E_s)$ over ten moves (pressure, spacing along the lateral, run time, nozzle, spacing between laterals, run time, pressure, layout, run time, spacing along the lateral); a move is accepted only if the score rises by more than 0.002. For pressure, nozzle and spacings the candidates are nine equally spaced values over the range of Table~\ref{tab:types}'s type; for the run time $T_0\times\{0.6,0.7,0.8,0.9,1,1.1,1.25\}$.

\begin{table}[!htbp]\centering\small\caption{First move of the optimiser: nine pressures from 200 to 500 kPa (rounded to 10) for the default impact layout; $\Delta$ is the change from the default score.}\label{tab:optP}
\begin{tabular}{lrrrrrr}\toprule
\hd{$P$ (kPa)} & \hd{$R$} & \hd{CU (\%)} & \hd{$T$} & \hd{$E_a$ (\%)} & \hd{Score} & \hd{$\Delta$}\\\midrule
200 & 12.8 & 82.1 & 548 & 67.9 & 0.6792 & -0.1405\\
240 & 14.0 & 90.7 & 449 & 75.5 & 0.7547 & -0.0650\\
280 & 15.1 & 95.0 & 391 & 80.4 & 0.8044 & -0.0153\\
310 & 15.9 & 96.3 & 363 & 82.6 & 0.8260 & 0.0063\\
350 & 16.9 & 96.5 & 348 & 81.0 & 0.8096 & -0.0101\\
390 & 17.9 & 95.8 & 335 & 79.5 & 0.7951 & -0.0246\\
420 & 18.5 & 95.5 & 324 & 79.3 & 0.7925 & -0.0272\\
460 & 19.4 & 95.4 & 308 & 79.7 & 0.7968 & -0.0229\\
500 & 20.2 & 95.7 & 293 & 80.5 & 0.8055 & -0.0142\\\bottomrule\end{tabular}\end{table}

\exbox{First optimiser move}{Read the first move from Table~\ref{tab:optP}: which pressure does it choose, and is the gain accepted?}\sol  The default (300 kPa) scores 0.8197. The highest score in the table is 0.8260 at 310 kPa, a change of 0.0063. A move is accepted only if the gain exceeds 0.002; here the gain is 0.0063, so it is accepted. The reason is that raising the pressure widens the reach (and so the overlap) while it also raises the drift loss; the two nearly cancel in $E_a$ when $E_s\approx100\%$ and the score equals $E_a$. \finans


\begin{table}[!htbp]\centering\small\caption{Sensitivity of $E_a$ (base 82.0\%) to the eight inputs of the layer.}\label{tab:sens}
\begin{tabular}{lrrrrr}\toprule
\hd{Input} & \hd{low factor} & \hd{$E_a$ (\%)} & \hd{high factor} & \hd{$E_a$ (\%)} & \hd{change (pp)}\\\midrule
Intake $k$ & 0.80 & 82.0 & 1.20 & 82.0 & 0.0\\
Wind & 0.50 & 83.0 & 1.50 & 80.8 & -2.2\\
Pressure & 0.80 & 75.5 & 1.20 & 80.5 & 5.1\\
Nozzle & 0.85 & 77.3 & 1.15 & 81.6 & 4.2\\
$S_s$ & 0.80 & 83.7 & 1.20 & 79.9 & -3.7\\
$S_l$ & 0.80 & 80.2 & 1.20 & 69.3 & -10.9\\
$Z_n$ & 0.75 & 79.6 & 1.25 & 83.5 & 3.9\\
Slope (via $S_d$) & 0.50 & 82.0 & 1.50 & 82.0 & 0.0\\\bottomrule\end{tabular}\end{table}

\begin{keybox}
In Table~\ref{tab:sens} the soil intake and the slope change nothing, because on the loam no point ponds, so runoff is zero in every case; the drivers are the spacings, the nozzle and the pressure (which change uniformity and the run time) and the depth $Z_n$ (a larger depth dilutes the fixed canopy loss). On a soil that ponds, intake and slope become the dominant inputs.
\end{keybox}
\exbox{Sets and water for the season}{The season of the loam at trigger $-300$ cm needs $I_y=211.68$ mm of irrigation (Manuscript~4, Fig.~3 data: 6.123 Mg ha$^{-1}$). For the 5 ha farm above, find the number of irrigations, the sprinklers, the sets, the round time and the gross water.}\sol \givenline{$I_y=211.68$ mm}\ans $n_{\rm ev}=I_y/Z_n=211.68/30=7.06\approx7$ irrigations. Sprinklers: $n_{\rm sp}=\mathrm{round}(10^4\times5/(12\times18))=231$. Total flow $Q_{\rm tot}=231\times21.23/60=81.7$ L s$^{-1}$; with $Q_s=40$ L s$^{-1}$ the farm needs $n_{\rm set}=\lceil81.7/40\rceil=3$ sets. Round time $t_{\rm round}=3\times371/60=18.6$ h. Gross depth per irrigation $=36.5$ mm; gross per season $=7.06\times36.5=258$ mm, i.e.\ $10\times258=2{,}576$ m$^3$ ha$^{-1}$. The gross is larger than the net $I_y$ because only $E_a=82.0\%$ of the discharge reaches the root zone usefully. \finans


\section{Kinds of system}
The layer draws five systems to scale: solid-set (fixed laterals, $E_a$ typically 70--75\%), hand-move (70--75\%), centre pivot (65--70\%), linear-move (65--70\%) and travelling gun (65--70\%). A bar locates the design's $E_a$ against the range of the chosen system. The default impact layout reaches 82\% on the loam in the model, which lies above the ranges quoted for fixed systems: the model counts no leakage, no pipe-end effect and no pressure variation along the lateral, so the figure is an upper estimate for a well-managed set.

\section{Summary of the sprinkler layer}
\begin{itemize}
\item Discharge $q\propto d^2\sqrt P$ (13.2); reach $R=c_R\sqrt{dh}f_\alpha$ (13.3).
\item Single-sprinkler pattern from $\hat p(x)$ with a wind stretch (13.4--13.5), overlapped on the grid and scored by CU (13.6--13.7); AR $=60q/(S_sS_l)$ exactly.
\item Drift $w_d$ (13.8), ponding time $t_p=(a_e/K)^{1/(A-1)}$ (13.9), canopy retention (13.10), a closed partition $E_a+\mathrm{DP}+\mathrm{RO}+\mathrm{WD}=1$ (13.11).
\item Run time by bisection (13.12), allowed rate from Table 14.3, ten checks (13.13--13.16), Blasius pipes and pump (13.17--13.20).
\end{itemize}
\IfFileExists{extra/app13.tex}{\subsection{The output of the application, panel by panel}
The sprinkler layer is organised in five parts: the site, the sprinkler and its layout, the way the system is run, the farm and its pipes, and the results with the course checks. The case shown is a solid-set system of impact sprinklers with 4.4 mm nozzles at 300 kPa, spaced 12 m along laterals 18 m apart, on 5 ha of coarse loam with a wind of 3 m\,s$^{-1}$; it applies 72 mm per irrigation.

The first part sets the conditions that limit the sprinkler. The soil tiles give the maximum application rate of each texture at the slope chosen (25 mm\,h$^{-1}$ for coarse loam on a 1.0 per cent slope), which is the rate the sprinkler must not exceed. The wind speed (3.0 m\,s$^{-1}$) enters the drift loss, and the surface storage (5.0 mm) is the depth of water the soil surface can hold before it runs.

\appfig{spr_site}{0.96}{Part A, your site: crop, mulch, soil type with the largest application rate it accepts for the slope, ground slope, wind and surface storage.}

The buttons choose the system (solid-set, hand-move, centre pivot, linear move, travelling gun) and the sprinkler (impact, rotating spray, micro-sprinkler, big gun) and the square or triangular layout. The numerical boxes set the nozzle diameter (4.4 mm), the trajectory angle (25$^\circ$), the pressure (300 kPa) and the spacing along and between laterals. The live view on the left draws the wetted circle (throw radius 15.7 m) and the one on the right the head of the sprinkler with the jet and the droplets; the text under it gives the flow (21.2 L\,min$^{-1}$), the jet speed (14 m\,s$^{-1}$), the biggest drop (3.8 mm) and the loss to wind and air (3.8 per cent).

\appfig{spr_layout}{0.96}{Part B, the sprinkler and its layout: system type, sprinkler type, layout, nozzle, trajectory angle, pressure and spacing, with the live view of the throw.}

Part C sets the depth per irrigation (72 mm, from the DSS) and the operating hours per day (16 h); the run time of 805 min is computed. Part D collects the area (5.00 ha), the lateral and mainline lengths, the farm supply (40 L\,s$^{-1}$) and the lift (12 m) from which the pipes and the pump are sized.

\appfig{spr_run}{0.96}{Parts C and D, how you run it and the farm with its water supply and pipes.}

The plan shows the sprinklers on the grid and the circles of the spray patterns overlapping. The colours give the depth infiltrated at each point: green is the depth needed, red-brown is short, blue is more than needed and brown hatching is water running off. The overlap of the patterns is the origin of the Christiansen uniformity of this chapter

\appfig{spr_plan}{0.96}{Layout in plan: depth infiltrated so far, from the overlapping spray patterns of the sprinklers.}

The section along the lateral shows the water that arrives at each point (the applied depth, including the drift loss), the depth that infiltrates and the volume of the soil that receives it; below the real texture of the soil the coloured column is the water that infiltrated, the depth divided by the water the soil can hold. The profile of the day at the bottom gives the trigger and field capacity for reference.

\appfig{spr_lateral}{0.96}{Along the lateral: depth applied, depth infiltrated and where it ends up in the soil, with the soil water of the day.}

The curve is the infiltration rate of the soil with time (the derivative of the Kostiakov intake of this chapter) and the horizontal line is the application rate of the sprinklers (5.9 mm\,h$^{-1}$, or 5.7 mm\,h$^{-1}$ after the drift loss). The soil intake at 30 min (19.2 mm\,h$^{-1}$) is well above the application rate, so no runoff occurs; the table lists ponding, runoff, surface storage and pattern uniformity.

\appfig{spr_intake}{0.96}{Soil intake against the application rate: when the soil cannot keep up, water runs off.}

The cards and the bar split the water that leaves the sprinklers into the part stored in the roots (91 per cent), the deep percolation, the drift and evaporation and any runoff; they are the quantities $E_a$, $E_s$ and $CU$ of the chapter.

\appfig{spr_result}{0.96}{Result for this design: application efficiency 90.6 per cent, root zone refilled 99.8 per cent, Christiansen uniformity 96.2 per cent, runoff 0, deep percolation 5.6 per cent, drift and evaporation 3.8 per cent.}

The layer applies the rules of the course as ten checks: the spacing along the lateral and between laterals against the wetted radius, the wind limit, the application rate against the maximum for the soil and against the saturated conductivity, the Christiansen uniformity, the lateral head loss (at most 20 per cent of the working head), the mainline gradient, the time to irrigate the whole field within the interval and the flow within the farm supply. A green tick means the rule is met, with the actual value written beside it.

\appfig{spr_checks}{0.96}{Does your design pass? The course checks, ten of ten satisfied.}

The teaching part of the layer begins with the kinds of system (solid-set, hand-move, centre pivot, linear move and travelling gun), each with its plan and its range of application efficiency (70 to 75 per cent for solid-set, 65 to 70 per cent for the self-propelled types), and shows where this design falls on the scale.

\appfig{spr_kinds}{0.96}{How a sprinkler works: the kinds of sprinkler system and their efficiency ranges.}

The panel names the parts of an impact sprinkler (the swing arm, the spring, the nozzle, the bearing and the riser) and classifies sprinklers by their operating pressure: low pressure (6 to 20 m), medium pressure (20 to 45 m, which includes this design) and high pressure (45 to 100 m, the guns). It repeats the discharge law $q=C_d\,a\sqrt{2gh}$ with $q=k\,h^{0.5}$ and the throw of the jet.

\appfig{spr_parts}{0.96}{The sprinkler: parts and pressure class, with the head of a rotary impact sprinkler.}

The farthest throw is at about 28 to 30 degrees; low angles (8 degrees) go under foliage and drift less. The panel shows how small droplets fall near the sprinkler and large ones far away (0.5 to 4 mm).

\appfig{spr_jet}{0.96}{The jet: trajectory angle, throw radius and droplet size.}

}{}
\IfFileExists{extra/long13.tex}{\subsection{Complete design: a solid-set sprinkler system for potatoes}
\exbox{Design of a solid-set system for 10 ha}{Ten hectares of potato are on the sandy loam S05. The root zone holds 12 mm net between irrigations, the wind is 3 m\,s$^{-1}$ and the air temperature 24$^\circ$C. Compare three sprinkler layouts, choose one, and find the number of sprinklers, the flow, the sets, the pressure and the time.}\sol
\ans Soil S05: $K_s=44.7$ mm\,h$^{-1}$, $\psi_f=49.6$ mm; Kostiakov fit $k=0.00310$, $A=0.778$. The soil can take water at the early rate and then slows, so the application rate of the sprinklers must not exceed the late intake
\ans Layout impact $d=4.4$ mm, $P=300$ kPa, $12\times18$ m: nozzle discharge $q=21.2$ L\,min$^{-1}$, throw $R=15.7$ m, mean application rate $5.9$ mm\,h$^{-1}$, $CU=96.2$ per cent, drift and evaporation loss $3.8$ per cent, set time $T=134$ min, $E_a=90.6$ per cent, runoff 0.0 per cent
\ans Layout impact $d=4.4$ mm, $P=300$ kPa, $12\times12$ m: nozzle discharge $q=21.2$ L\,min$^{-1}$, throw $R=15.7$ m, mean application rate $8.8$ mm\,h$^{-1}$, $CU=98.9$ per cent, drift and evaporation loss $3.8$ per cent, set time $T=86$ min, $E_a=94.7$ per cent, runoff 0.0 per cent
\ans Layout spray $d=3.2$ mm, $P=170$ kPa, $6\times6$ m: nozzle discharge $q=8.5$ L\,min$^{-1}$, throw $R=6.0$ m, mean application rate $14.1$ mm\,h$^{-1}$, $CU=78.1$ per cent, drift and evaporation loss $4.1$ per cent, set time $T=75$ min, $E_a=67.1$ per cent, runoff 0.0 per cent
\ans Choice: impact at $12\times12$ m, giving $E_a=94.7$ per cent
\ans Sprinklers per 10 ha: $100\,000/(12\times12)=694$; if all operate together the flow is $694\times21.2/60=245.7$ L\,s$^{-1}$
\ans Divide the field into 4 sets, one at a time: flow $61.4$ L\,s$^{-1}$, a pump of this capacity at the sprinkler pressure $300$ kPa; add 20 per cent for pipe losses, $360$ kPa at the pump
\ans Set time $T=86$ min; $4$ sets $\times$ $86$ min $=5.7$ h per irrigation, so the system works 5.7 h per interval of 3 days
\ans Gross depth per application $=\overline{AR}\,T=12.7$ mm, of which 12.0 mm is stored and $0.2$ mm lost as deep percolation


}{}
\IfFileExists{extra/proj13.tex}{\banner{Design project 13: a solid-set sprinkler block for alfalfa}\label{sec:proj13}
\textbf{Brief.} A 12-ha block of alfalfa on a sandy loam (basic intake 15 mm\,h$^{-1}$) is to be irrigated by a solid-set system with sprinklers on a $12\times18$ m rectangle. The peak $ET_c$ is 8 mm\,d$^{-1}$ and the application efficiency is 0.80, so the gross depth for a 5-day interval is $8\times5/0.8=50$ mm. The designer compares four nozzles for the application rate.

\begin{table}[!htbp]\centering\small\caption{Sprinkler nozzle alternatives on a $12\times18$ m rectangular spacing.}
\begin{tabular}{rrrrrr}\toprule
\hd{Nozzle (mm)} & \hd{Pressure (kPa)} & \hd{$q$ (m$^3$/h)} & \hd{$R$ (m)} & \hd{App. rate (mm/h)} & \hd{$S_l/2R$}\\\midrule
4.0 & 300 & 1.05 & 11.8 & 4.9 & 0.51\\
4.8 & 350 & 1.64 & 13.9 & 7.6 & 0.43\\
5.6 & 350 & 2.23 & 15.1 & 10.3 & 0.40\\
6.4 & 400 & 3.11 & 17.2 & 14.4 & 0.35\\\bottomrule\end{tabular}\end{table}

\textbf{Reading.} The nozzle of 4.8 mm at 350 kPa gives 7.6 mm\,h$^{-1}$, below the intake rate of 15 mm\,h$^{-1}$, so there is no runoff; the set time for the 50 mm is $6.6$ h, and two sets a day complete the block. The number of sprinklers in operation is the block area divided by the spacing, $12\times10^4/216=555$, for a system flow of 910 m$^3$\,h$^{-1}$ if the whole block runs at once.
\fig{g13_pat}{0.7}{Radial application patterns of three sprinkler types, scaled to a wetted radius of 14 m.}
}{}
\IfFileExists{extra/ten13.tex}{\subsection{Worked solutions on the ten soils}
The same nozzle and spacing are tested on each soil.

\exbox{Ten-soil example 13a: Coarse loam (S01)}{Coarse loam ($K_s=13.0$ mm\,h$^{-1}$): a 4.8-mm nozzle at 350 kPa works on a $12\times18$ m spacing. Find the discharge, the application rate and whether runoff will occur.}\sol $h=350/9.81=35.7$ m; $q=0.95\pi(35.7/2000)^2\sqrt{2\times9.81\times4.8}\times6\times10^4=27.3$ L\,min$^{-1}$ \ans $r=q\times60/(S_lS_m)=27.3\times60/216=7.6$ mm\,h$^{-1}$ \ans The rate is below $K_s$, so no runoff. \ans

\exbox{Ten-soil example 13a: Coarse sand (S02)}{Coarse sand ($K_s=120.0$ mm\,h$^{-1}$): a 4.8-mm nozzle at 350 kPa works on a $12\times18$ m spacing. Find the discharge, the application rate and whether runoff will occur.}\sol $h=350/9.81=35.7$ m; $q=0.95\pi(35.7/2000)^2\sqrt{2\times9.81\times4.8}\times6\times10^4=27.3$ L\,min$^{-1}$ \ans $r=q\times60/(S_lS_m)=27.3\times60/216=7.6$ mm\,h$^{-1}$ \ans The rate is below $K_s$, so no runoff. \ans

\exbox{Ten-soil example 13a: Sand (S03)}{Sand ($K_s=210.0$ mm\,h$^{-1}$): a 4.8-mm nozzle at 350 kPa works on a $12\times18$ m spacing. Find the discharge, the application rate and whether runoff will occur.}\sol $h=350/9.81=35.7$ m; $q=0.95\pi(35.7/2000)^2\sqrt{2\times9.81\times4.8}\times6\times10^4=27.3$ L\,min$^{-1}$ \ans $r=q\times60/(S_lS_m)=27.3\times60/216=7.6$ mm\,h$^{-1}$ \ans The rate is below $K_s$, so no runoff. \ans

\exbox{Ten-soil example 13a: Loamy sand (S04)}{Loamy sand ($K_s=61.0$ mm\,h$^{-1}$): a 4.8-mm nozzle at 350 kPa works on a $12\times18$ m spacing. Find the discharge, the application rate and whether runoff will occur.}\sol $h=350/9.81=35.7$ m; $q=0.95\pi(35.7/2000)^2\sqrt{2\times9.81\times4.8}\times6\times10^4=27.3$ L\,min$^{-1}$ \ans $r=q\times60/(S_lS_m)=27.3\times60/216=7.6$ mm\,h$^{-1}$ \ans The rate is below $K_s$, so no runoff. \ans

\exbox{Ten-soil example 13a: Sandy loam (S05)}{Sandy loam ($K_s=26.0$ mm\,h$^{-1}$): a 4.8-mm nozzle at 350 kPa works on a $12\times18$ m spacing. Find the discharge, the application rate and whether runoff will occur.}\sol $h=350/9.81=35.7$ m; $q=0.95\pi(35.7/2000)^2\sqrt{2\times9.81\times4.8}\times6\times10^4=27.3$ L\,min$^{-1}$ \ans $r=q\times60/(S_lS_m)=27.3\times60/216=7.6$ mm\,h$^{-1}$ \ans The rate is below $K_s$, so no runoff. \ans

\exbox{Ten-soil example 13a: Loam (S06)}{Loam ($K_s=13.0$ mm\,h$^{-1}$): a 4.8-mm nozzle at 350 kPa works on a $12\times18$ m spacing. Find the discharge, the application rate and whether runoff will occur.}\sol $h=350/9.81=35.7$ m; $q=0.95\pi(35.7/2000)^2\sqrt{2\times9.81\times4.8}\times6\times10^4=27.3$ L\,min$^{-1}$ \ans $r=q\times60/(S_lS_m)=27.3\times60/216=7.6$ mm\,h$^{-1}$ \ans The rate is below $K_s$, so no runoff. \ans

\exbox{Ten-soil example 13a: Silt loam (S07)}{Silt loam ($K_s=6.8$ mm\,h$^{-1}$): a 4.8-mm nozzle at 350 kPa works on a $12\times18$ m spacing. Find the discharge, the application rate and whether runoff will occur.}\sol $h=350/9.81=35.7$ m; $q=0.95\pi(35.7/2000)^2\sqrt{2\times9.81\times4.8}\times6\times10^4=27.3$ L\,min$^{-1}$ \ans $r=q\times60/(S_lS_m)=27.3\times60/216=7.6$ mm\,h$^{-1}$ \ans The rate exceeds $K_s$: ponding after $68.4$ h; sets longer than that run off, so use a smaller nozzle or a wider spacing. \ans

\exbox{Ten-soil example 13a: Sandy clay loam (S08)}{Sandy clay loam ($K_s=4.3$ mm\,h$^{-1}$): a 4.8-mm nozzle at 350 kPa works on a $12\times18$ m spacing. Find the discharge, the application rate and whether runoff will occur.}\sol $h=350/9.81=35.7$ m; $q=0.95\pi(35.7/2000)^2\sqrt{2\times9.81\times4.8}\times6\times10^4=27.3$ L\,min$^{-1}$ \ans $r=q\times60/(S_lS_m)=27.3\times60/216=7.6$ mm\,h$^{-1}$ \ans The rate exceeds $K_s$: ponding after $11.4$ h; sets longer than that run off, so use a smaller nozzle or a wider spacing. \ans

\exbox{Ten-soil example 13a: Clay loam (S09)}{Clay loam ($K_s=2.3$ mm\,h$^{-1}$): a 4.8-mm nozzle at 350 kPa works on a $12\times18$ m spacing. Find the discharge, the application rate and whether runoff will occur.}\sol $h=350/9.81=35.7$ m; $q=0.95\pi(35.7/2000)^2\sqrt{2\times9.81\times4.8}\times6\times10^4=27.3$ L\,min$^{-1}$ \ans $r=q\times60/(S_lS_m)=27.3\times60/216=7.6$ mm\,h$^{-1}$ \ans The rate exceeds $K_s$: ponding after $3.7$ h; sets longer than that run off, so use a smaller nozzle or a wider spacing. \ans

\exbox{Ten-soil example 13a: Silty clay loam (S10)}{Silty clay loam ($K_s=1.5$ mm\,h$^{-1}$): a 4.8-mm nozzle at 350 kPa works on a $12\times18$ m spacing. Find the discharge, the application rate and whether runoff will occur.}\sol $h=350/9.81=35.7$ m; $q=0.95\pi(35.7/2000)^2\sqrt{2\times9.81\times4.8}\times6\times10^4=27.3$ L\,min$^{-1}$ \ans $r=q\times60/(S_lS_m)=27.3\times60/216=7.6$ mm\,h$^{-1}$ \ans The rate exceeds $K_s$: ponding after $2.7$ h; sets longer than that run off, so use a smaller nozzle or a wider spacing. \ans

}{}
\IfFileExists{extra/tenb13.tex}{\subsection{Worked solutions on ten crops}
The ten crops of the library differ in their demand, their roots and their sensitivity. The same type of question is answered for each of them at $ET_0=5.5$ mm\,d$^{-1}$; the data are typical values of the crops (Doorenbos--Kassam and FAO 56).

\exbox{Crop example 13a: Maize}{Maize: a solid-set sprinkler system with $r=7.6$ mm\,h$^{-1}$ irrigates every 3 days at $E_a=0.75$. Find the gross depth, the set time and the number of sets per 20-h day.}\sol \givenline{$r=7.6$ mm\,h$^{-1}$ irrigates every ; $E_a=0.75$}\ans Net depth $3\times6.60=19.8$ mm; gross $19.8/0.75=26.4$ mm \ans $t_{\rm set}=26.4/7.6=3.5$ h \ans Sets per day $\lfloor20/3.5\rfloor=5$ \finans

\exbox{Crop example 13a: Wheat}{Wheat: a solid-set sprinkler system with $r=7.6$ mm\,h$^{-1}$ irrigates every 3 days at $E_a=0.75$. Find the gross depth, the set time and the number of sets per 20-h day.}\sol \givenline{$r=7.6$ mm\,h$^{-1}$ irrigates every ; $E_a=0.75$}\ans Net depth $3\times6.32=19.0$ mm; gross $19.0/0.75=25.3$ mm \ans $t_{\rm set}=25.3/7.6=3.3$ h \ans Sets per day $\lfloor20/3.3\rfloor=6$ \finans

\exbox{Crop example 13a: Cotton}{Cotton: a solid-set sprinkler system with $r=7.6$ mm\,h$^{-1}$ irrigates every 3 days at $E_a=0.75$. Find the gross depth, the set time and the number of sets per 20-h day.}\sol \givenline{$r=7.6$ mm\,h$^{-1}$ irrigates every ; $E_a=0.75$}\ans Net depth $3\times6.32=19.0$ mm; gross $19.0/0.75=25.3$ mm \ans $t_{\rm set}=25.3/7.6=3.3$ h \ans Sets per day $\lfloor20/3.3\rfloor=6$ \finans

\exbox{Crop example 13a: Tomato}{Tomato: a solid-set sprinkler system with $r=7.6$ mm\,h$^{-1}$ irrigates every 3 days at $E_a=0.75$. Find the gross depth, the set time and the number of sets per 20-h day.}\sol \givenline{$r=7.6$ mm\,h$^{-1}$ irrigates every ; $E_a=0.75$}\ans Net depth $3\times6.32=19.0$ mm; gross $19.0/0.75=25.3$ mm \ans $t_{\rm set}=25.3/7.6=3.3$ h \ans Sets per day $\lfloor20/3.3\rfloor=6$ \finans

\exbox{Crop example 13a: Potato}{Potato: a solid-set sprinkler system with $r=7.6$ mm\,h$^{-1}$ irrigates every 3 days at $E_a=0.75$. Find the gross depth, the set time and the number of sets per 20-h day.}\sol \givenline{$r=7.6$ mm\,h$^{-1}$ irrigates every ; $E_a=0.75$}\ans Net depth $3\times6.32=19.0$ mm; gross $19.0/0.75=25.3$ mm \ans $t_{\rm set}=25.3/7.6=3.3$ h \ans Sets per day $\lfloor20/3.3\rfloor=6$ \finans

\exbox{Crop example 13a: Soybean}{Soybean: a solid-set sprinkler system with $r=7.6$ mm\,h$^{-1}$ irrigates every 3 days at $E_a=0.75$. Find the gross depth, the set time and the number of sets per 20-h day.}\sol \givenline{$r=7.6$ mm\,h$^{-1}$ irrigates every ; $E_a=0.75$}\ans Net depth $3\times6.32=19.0$ mm; gross $19.0/0.75=25.3$ mm \ans $t_{\rm set}=25.3/7.6=3.3$ h \ans Sets per day $\lfloor20/3.3\rfloor=6$ \finans

\exbox{Crop example 13a: Sugar beet}{Sugar beet: a solid-set sprinkler system with $r=7.6$ mm\,h$^{-1}$ irrigates every 3 days at $E_a=0.75$. Find the gross depth, the set time and the number of sets per 20-h day.}\sol \givenline{$r=7.6$ mm\,h$^{-1}$ irrigates every ; $E_a=0.75$}\ans Net depth $3\times6.60=19.8$ mm; gross $19.8/0.75=26.4$ mm \ans $t_{\rm set}=26.4/7.6=3.5$ h \ans Sets per day $\lfloor20/3.5\rfloor=5$ \finans

\exbox{Crop example 13a: Sunflower}{Sunflower: a solid-set sprinkler system with $r=7.6$ mm\,h$^{-1}$ irrigates every 3 days at $E_a=0.75$. Find the gross depth, the set time and the number of sets per 20-h day.}\sol \givenline{$r=7.6$ mm\,h$^{-1}$ irrigates every ; $E_a=0.75$}\ans Net depth $3\times6.05=18.2$ mm; gross $18.2/0.75=24.2$ mm \ans $t_{\rm set}=24.2/7.6=3.2$ h \ans Sets per day $\lfloor20/3.2\rfloor=6$ \finans

\exbox{Crop example 13a: Onion}{Onion: a solid-set sprinkler system with $r=7.6$ mm\,h$^{-1}$ irrigates every 3 days at $E_a=0.75$. Find the gross depth, the set time and the number of sets per 20-h day.}\sol \givenline{$r=7.6$ mm\,h$^{-1}$ irrigates every ; $E_a=0.75$}\ans Net depth $3\times5.78=17.3$ mm; gross $17.3/0.75=23.1$ mm \ans $t_{\rm set}=23.1/7.6=3.0$ h \ans Sets per day $\lfloor20/3.0\rfloor=6$ \finans

\exbox{Crop example 13a: Alfalfa}{Alfalfa: a solid-set sprinkler system with $r=7.6$ mm\,h$^{-1}$ irrigates every 3 days at $E_a=0.75$. Find the gross depth, the set time and the number of sets per 20-h day.}\sol \givenline{$r=7.6$ mm\,h$^{-1}$ irrigates every ; $E_a=0.75$}\ans Net depth $3\times6.32=19.0$ mm; gross $19.0/0.75=25.3$ mm \ans $t_{\rm set}=25.3/7.6=3.3$ h \ans Sets per day $\lfloor20/3.3\rfloor=6$ \finans

}{}
\IfFileExists{extra/teny13.tex}{\subsection{Worked solutions on the ten soils (third scenario)}
The same questions are repeated for different working conditions, to show how the answer moves with them.

\exbox{Ten-soil example 13c: Coarse loam (S01)}{Coarse loam ($K_s=13.0$ mm\,h$^{-1}$): a 5.6-mm nozzle at 400 kPa works on a $12\times18$ m spacing. Find the discharge, the application rate and whether runoff will occur.}\sol \givenline{$K_s=13.0$ mm\,h$^{-1}$}\ans $h=350/9.81=40.8$ m; \[q=0.95\pi(40.8/2000)^2\sqrt{2\times9.81\times5.6}\times6\times10^4=39.7\ \text{L\,min$^{-1}$}\] \ans \[r=q\times60/(S_lS_m)=39.7\times60/216=11.0\ \text{mm\,h$^{-1}$}\] \ans The rate is below $K_s$, so no runoff. \finans

\exbox{Ten-soil example 13c: Coarse sand (S02)}{Coarse sand ($K_s=120.0$ mm\,h$^{-1}$): a 5.6-mm nozzle at 400 kPa works on a $12\times18$ m spacing. Find the discharge, the application rate and whether runoff will occur.}\sol \givenline{$K_s=120.0$ mm\,h$^{-1}$}\ans $h=350/9.81=40.8$ m; \[q=0.95\pi(40.8/2000)^2\sqrt{2\times9.81\times5.6}\times6\times10^4=39.7\ \text{L\,min$^{-1}$}\] \ans \[r=q\times60/(S_lS_m)=39.7\times60/216=11.0\ \text{mm\,h$^{-1}$}\] \ans The rate is below $K_s$, so no runoff. \finans

\exbox{Ten-soil example 13c: Sand (S03)}{Sand ($K_s=210.0$ mm\,h$^{-1}$): a 5.6-mm nozzle at 400 kPa works on a $12\times18$ m spacing. Find the discharge, the application rate and whether runoff will occur.}\sol \givenline{$K_s=210.0$ mm\,h$^{-1}$}\ans $h=350/9.81=40.8$ m; \[q=0.95\pi(40.8/2000)^2\sqrt{2\times9.81\times5.6}\times6\times10^4=39.7\ \text{L\,min$^{-1}$}\] \ans \[r=q\times60/(S_lS_m)=39.7\times60/216=11.0\ \text{mm\,h$^{-1}$}\] \ans The rate is below $K_s$, so no runoff. \finans

\exbox{Ten-soil example 13c: Loamy sand (S04)}{Loamy sand ($K_s=61.0$ mm\,h$^{-1}$): a 5.6-mm nozzle at 400 kPa works on a $12\times18$ m spacing. Find the discharge, the application rate and whether runoff will occur.}\sol \givenline{$K_s=61.0$ mm\,h$^{-1}$}\ans $h=350/9.81=40.8$ m; \[q=0.95\pi(40.8/2000)^2\sqrt{2\times9.81\times5.6}\times6\times10^4=39.7\ \text{L\,min$^{-1}$}\] \ans \[r=q\times60/(S_lS_m)=39.7\times60/216=11.0\ \text{mm\,h$^{-1}$}\] \ans The rate is below $K_s$, so no runoff. \finans

\exbox{Ten-soil example 13c: Sandy loam (S05)}{Sandy loam ($K_s=26.0$ mm\,h$^{-1}$): a 5.6-mm nozzle at 400 kPa works on a $12\times18$ m spacing. Find the discharge, the application rate and whether runoff will occur.}\sol \givenline{$K_s=26.0$ mm\,h$^{-1}$}\ans $h=350/9.81=40.8$ m; \[q=0.95\pi(40.8/2000)^2\sqrt{2\times9.81\times5.6}\times6\times10^4=39.7\ \text{L\,min$^{-1}$}\] \ans \[r=q\times60/(S_lS_m)=39.7\times60/216=11.0\ \text{mm\,h$^{-1}$}\] \ans The rate is below $K_s$, so no runoff. \finans

\exbox{Ten-soil example 13c: Loam (S06)}{Loam ($K_s=13.0$ mm\,h$^{-1}$): a 5.6-mm nozzle at 400 kPa works on a $12\times18$ m spacing. Find the discharge, the application rate and whether runoff will occur.}\sol \givenline{$K_s=13.0$ mm\,h$^{-1}$}\ans $h=350/9.81=40.8$ m; \[q=0.95\pi(40.8/2000)^2\sqrt{2\times9.81\times5.6}\times6\times10^4=39.7\ \text{L\,min$^{-1}$}\] \ans \[r=q\times60/(S_lS_m)=39.7\times60/216=11.0\ \text{mm\,h$^{-1}$}\] \ans The rate is below $K_s$, so no runoff. \finans

\exbox{Ten-soil example 13c: Silt loam (S07)}{Silt loam ($K_s=6.8$ mm\,h$^{-1}$): a 5.6-mm nozzle at 400 kPa works on a $12\times18$ m spacing. Find the discharge, the application rate and whether runoff will occur.}\sol \givenline{$K_s=6.8$ mm\,h$^{-1}$}\ans $h=350/9.81=40.8$ m; \[q=0.95\pi(40.8/2000)^2\sqrt{2\times9.81\times5.6}\times6\times10^4=39.7\ \text{L\,min$^{-1}$}\] \ans \[r=q\times60/(S_lS_m)=39.7\times60/216=11.0\ \text{mm\,h$^{-1}$}\] \ans The rate exceeds $K_s$: ponding after $8.7$ h; sets longer than that run off, so use a smaller nozzle or a wider spacing. \finans

\exbox{Ten-soil example 13c: Sandy clay loam (S08)}{Sandy clay loam ($K_s=4.3$ mm\,h$^{-1}$): a 5.6-mm nozzle at 400 kPa works on a $12\times18$ m spacing. Find the discharge, the application rate and whether runoff will occur.}\sol \givenline{$K_s=4.3$ mm\,h$^{-1}$}\ans $h=350/9.81=40.8$ m; \[q=0.95\pi(40.8/2000)^2\sqrt{2\times9.81\times5.6}\times6\times10^4=39.7\ \text{L\,min$^{-1}$}\] \ans \[r=q\times60/(S_lS_m)=39.7\times60/216=11.0\ \text{mm\,h$^{-1}$}\] \ans The rate exceeds $K_s$: ponding after $3.8$ h; sets longer than that run off, so use a smaller nozzle or a wider spacing. \finans

\exbox{Ten-soil example 13c: Clay loam (S09)}{Clay loam ($K_s=2.3$ mm\,h$^{-1}$): a 5.6-mm nozzle at 400 kPa works on a $12\times18$ m spacing. Find the discharge, the application rate and whether runoff will occur.}\sol \givenline{$K_s=2.3$ mm\,h$^{-1}$}\ans $h=350/9.81=40.8$ m; \[q=0.95\pi(40.8/2000)^2\sqrt{2\times9.81\times5.6}\times6\times10^4=39.7\ \text{L\,min$^{-1}$}\] \ans \[r=q\times60/(S_lS_m)=39.7\times60/216=11.0\ \text{mm\,h$^{-1}$}\] \ans The rate exceeds $K_s$: ponding after $1.5$ h; sets longer than that run off, so use a smaller nozzle or a wider spacing. \finans

\exbox{Ten-soil example 13c: Silty clay loam (S10)}{Silty clay loam ($K_s=1.5$ mm\,h$^{-1}$): a 5.6-mm nozzle at 400 kPa works on a $12\times18$ m spacing. Find the discharge, the application rate and whether runoff will occur.}\sol \givenline{$K_s=1.5$ mm\,h$^{-1}$}\ans $h=350/9.81=40.8$ m; \[q=0.95\pi(40.8/2000)^2\sqrt{2\times9.81\times5.6}\times6\times10^4=39.7\ \text{L\,min$^{-1}$}\] \ans \[r=q\times60/(S_lS_m)=39.7\times60/216=11.0\ \text{mm\,h$^{-1}$}\] \ans The rate exceeds $K_s$: ponding after $1.2$ h; sets longer than that run off, so use a smaller nozzle or a wider spacing. \finans

}{}
\IfFileExists{extra/tenby13.tex}{\subsection{Worked solutions on ten crops (hot climate)}
The ten crops of the library differ in their demand, their roots and their sensitivity. The same type of question is answered for each of them at $ET_0=7.5$ mm\,d$^{-1}$; the data are typical values of the crops (Doorenbos--Kassam and FAO 56).

\exbox{Crop example 13c: Maize}{Maize: a solid-set sprinkler system with $r=7.6$ mm\,h$^{-1}$ irrigates every 3 days at $E_a=0.75$. Find the gross depth, the set time and the number of sets per 20-h day.}\sol Net depth $3\times9.00=27.0$ mm; gross $27.0/0.75=36.0$ mm \ans $t_{\rm set}=36.0/7.6=4.7$ h \ans Sets per day $\lfloor20/4.7\rfloor=4$ \ans

\exbox{Crop example 13c: Wheat}{Wheat: a solid-set sprinkler system with $r=7.6$ mm\,h$^{-1}$ irrigates every 3 days at $E_a=0.75$. Find the gross depth, the set time and the number of sets per 20-h day.}\sol Net depth $3\times8.62=25.9$ mm; gross $25.9/0.75=34.5$ mm \ans $t_{\rm set}=34.5/7.6=4.5$ h \ans Sets per day $\lfloor20/4.5\rfloor=4$ \ans

\exbox{Crop example 13c: Cotton}{Cotton: a solid-set sprinkler system with $r=7.6$ mm\,h$^{-1}$ irrigates every 3 days at $E_a=0.75$. Find the gross depth, the set time and the number of sets per 20-h day.}\sol Net depth $3\times8.62=25.9$ mm; gross $25.9/0.75=34.5$ mm \ans $t_{\rm set}=34.5/7.6=4.5$ h \ans Sets per day $\lfloor20/4.5\rfloor=4$ \ans

\exbox{Crop example 13c: Tomato}{Tomato: a solid-set sprinkler system with $r=7.6$ mm\,h$^{-1}$ irrigates every 3 days at $E_a=0.75$. Find the gross depth, the set time and the number of sets per 20-h day.}\sol Net depth $3\times8.62=25.9$ mm; gross $25.9/0.75=34.5$ mm \ans $t_{\rm set}=34.5/7.6=4.5$ h \ans Sets per day $\lfloor20/4.5\rfloor=4$ \ans

\exbox{Crop example 13c: Potato}{Potato: a solid-set sprinkler system with $r=7.6$ mm\,h$^{-1}$ irrigates every 3 days at $E_a=0.75$. Find the gross depth, the set time and the number of sets per 20-h day.}\sol Net depth $3\times8.62=25.9$ mm; gross $25.9/0.75=34.5$ mm \ans $t_{\rm set}=34.5/7.6=4.5$ h \ans Sets per day $\lfloor20/4.5\rfloor=4$ \ans

\exbox{Crop example 13c: Soybean}{Soybean: a solid-set sprinkler system with $r=7.6$ mm\,h$^{-1}$ irrigates every 3 days at $E_a=0.75$. Find the gross depth, the set time and the number of sets per 20-h day.}\sol Net depth $3\times8.62=25.9$ mm; gross $25.9/0.75=34.5$ mm \ans $t_{\rm set}=34.5/7.6=4.5$ h \ans Sets per day $\lfloor20/4.5\rfloor=4$ \ans

\exbox{Crop example 13c: Sugar beet}{Sugar beet: a solid-set sprinkler system with $r=7.6$ mm\,h$^{-1}$ irrigates every 3 days at $E_a=0.75$. Find the gross depth, the set time and the number of sets per 20-h day.}\sol Net depth $3\times9.00=27.0$ mm; gross $27.0/0.75=36.0$ mm \ans $t_{\rm set}=36.0/7.6=4.7$ h \ans Sets per day $\lfloor20/4.7\rfloor=4$ \ans

\exbox{Crop example 13c: Sunflower}{Sunflower: a solid-set sprinkler system with $r=7.6$ mm\,h$^{-1}$ irrigates every 3 days at $E_a=0.75$. Find the gross depth, the set time and the number of sets per 20-h day.}\sol Net depth $3\times8.25=24.8$ mm; gross $24.8/0.75=33.0$ mm \ans $t_{\rm set}=33.0/7.6=4.3$ h \ans Sets per day $\lfloor20/4.3\rfloor=4$ \ans

\exbox{Crop example 13c: Onion}{Onion: a solid-set sprinkler system with $r=7.6$ mm\,h$^{-1}$ irrigates every 3 days at $E_a=0.75$. Find the gross depth, the set time and the number of sets per 20-h day.}\sol Net depth $3\times7.88=23.6$ mm; gross $23.6/0.75=31.5$ mm \ans $t_{\rm set}=31.5/7.6=4.1$ h \ans Sets per day $\lfloor20/4.1\rfloor=4$ \ans

\exbox{Crop example 13c: Alfalfa}{Alfalfa: a solid-set sprinkler system with $r=7.6$ mm\,h$^{-1}$ irrigates every 3 days at $E_a=0.75$. Find the gross depth, the set time and the number of sets per 20-h day.}\sol Net depth $3\times8.62=25.9$ mm; gross $25.9/0.75=34.5$ mm \ans $t_{\rm set}=34.5/7.6=4.5$ h \ans Sets per day $\lfloor20/4.5\rfloor=4$ \ans

}{}
\IfFileExists{extra/tenx13.tex}{\section{Worked solutions on the ten soils (second scenario)}
The same questions are repeated for different working conditions, to show how the answer moves with them.

\exbox{Ten-soil example 13b: Coarse loam (S01)}{Coarse loam ($K_s=13.0$ mm\,h$^{-1}$): a 4.0-mm nozzle at 300 kPa works on a $12\times18$ m spacing. Find the discharge, the application rate and whether runoff will occur.}\sol $h=350/9.81=30.6$ m; $q=0.95\pi(30.6/2000)^2\sqrt{2\times9.81\times4.0}\times6\times10^4=17.5$ L\,min$^{-1}$ \ans $r=q\times60/(S_lS_m)=17.5\times60/216=4.9$ mm\,h$^{-1}$ \ans The rate is below $K_s$, so no runoff. \ans

\exbox{Ten-soil example 13b: Coarse sand (S02)}{Coarse sand ($K_s=120.0$ mm\,h$^{-1}$): a 4.0-mm nozzle at 300 kPa works on a $12\times18$ m spacing. Find the discharge, the application rate and whether runoff will occur.}\sol $h=350/9.81=30.6$ m; $q=0.95\pi(30.6/2000)^2\sqrt{2\times9.81\times4.0}\times6\times10^4=17.5$ L\,min$^{-1}$ \ans $r=q\times60/(S_lS_m)=17.5\times60/216=4.9$ mm\,h$^{-1}$ \ans The rate is below $K_s$, so no runoff. \ans

\exbox{Ten-soil example 13b: Sand (S03)}{Sand ($K_s=210.0$ mm\,h$^{-1}$): a 4.0-mm nozzle at 300 kPa works on a $12\times18$ m spacing. Find the discharge, the application rate and whether runoff will occur.}\sol $h=350/9.81=30.6$ m; $q=0.95\pi(30.6/2000)^2\sqrt{2\times9.81\times4.0}\times6\times10^4=17.5$ L\,min$^{-1}$ \ans $r=q\times60/(S_lS_m)=17.5\times60/216=4.9$ mm\,h$^{-1}$ \ans The rate is below $K_s$, so no runoff. \ans

\exbox{Ten-soil example 13b: Loamy sand (S04)}{Loamy sand ($K_s=61.0$ mm\,h$^{-1}$): a 4.0-mm nozzle at 300 kPa works on a $12\times18$ m spacing. Find the discharge, the application rate and whether runoff will occur.}\sol $h=350/9.81=30.6$ m; $q=0.95\pi(30.6/2000)^2\sqrt{2\times9.81\times4.0}\times6\times10^4=17.5$ L\,min$^{-1}$ \ans $r=q\times60/(S_lS_m)=17.5\times60/216=4.9$ mm\,h$^{-1}$ \ans The rate is below $K_s$, so no runoff. \ans

\exbox{Ten-soil example 13b: Sandy loam (S05)}{Sandy loam ($K_s=26.0$ mm\,h$^{-1}$): a 4.0-mm nozzle at 300 kPa works on a $12\times18$ m spacing. Find the discharge, the application rate and whether runoff will occur.}\sol $h=350/9.81=30.6$ m; $q=0.95\pi(30.6/2000)^2\sqrt{2\times9.81\times4.0}\times6\times10^4=17.5$ L\,min$^{-1}$ \ans $r=q\times60/(S_lS_m)=17.5\times60/216=4.9$ mm\,h$^{-1}$ \ans The rate is below $K_s$, so no runoff. \ans

\exbox{Ten-soil example 13b: Loam (S06)}{Loam ($K_s=13.0$ mm\,h$^{-1}$): a 4.0-mm nozzle at 300 kPa works on a $12\times18$ m spacing. Find the discharge, the application rate and whether runoff will occur.}\sol $h=350/9.81=30.6$ m; $q=0.95\pi(30.6/2000)^2\sqrt{2\times9.81\times4.0}\times6\times10^4=17.5$ L\,min$^{-1}$ \ans $r=q\times60/(S_lS_m)=17.5\times60/216=4.9$ mm\,h$^{-1}$ \ans The rate is below $K_s$, so no runoff. \ans

\exbox{Ten-soil example 13b: Silt loam (S07)}{Silt loam ($K_s=6.8$ mm\,h$^{-1}$): a 4.0-mm nozzle at 300 kPa works on a $12\times18$ m spacing. Find the discharge, the application rate and whether runoff will occur.}\sol $h=350/9.81=30.6$ m; $q=0.95\pi(30.6/2000)^2\sqrt{2\times9.81\times4.0}\times6\times10^4=17.5$ L\,min$^{-1}$ \ans $r=q\times60/(S_lS_m)=17.5\times60/216=4.9$ mm\,h$^{-1}$ \ans The rate is below $K_s$, so no runoff. \ans

\exbox{Ten-soil example 13b: Sandy clay loam (S08)}{Sandy clay loam ($K_s=4.3$ mm\,h$^{-1}$): a 4.0-mm nozzle at 300 kPa works on a $12\times18$ m spacing. Find the discharge, the application rate and whether runoff will occur.}\sol $h=350/9.81=30.6$ m; $q=0.95\pi(30.6/2000)^2\sqrt{2\times9.81\times4.0}\times6\times10^4=17.5$ L\,min$^{-1}$ \ans $r=q\times60/(S_lS_m)=17.5\times60/216=4.9$ mm\,h$^{-1}$ \ans The rate exceeds $K_s$: ponding after $101.0$ h; sets longer than that run off, so use a smaller nozzle or a wider spacing. \ans

\exbox{Ten-soil example 13b: Clay loam (S09)}{Clay loam ($K_s=2.3$ mm\,h$^{-1}$): a 4.0-mm nozzle at 300 kPa works on a $12\times18$ m spacing. Find the discharge, the application rate and whether runoff will occur.}\sol $h=350/9.81=30.6$ m; $q=0.95\pi(30.6/2000)^2\sqrt{2\times9.81\times4.0}\times6\times10^4=17.5$ L\,min$^{-1}$ \ans $r=q\times60/(S_lS_m)=17.5\times60/216=4.9$ mm\,h$^{-1}$ \ans The rate exceeds $K_s$: ponding after $11.9$ h; sets longer than that run off, so use a smaller nozzle or a wider spacing. \ans

\exbox{Ten-soil example 13b: Silty clay loam (S10)}{Silty clay loam ($K_s=1.5$ mm\,h$^{-1}$): a 4.0-mm nozzle at 300 kPa works on a $12\times18$ m spacing. Find the discharge, the application rate and whether runoff will occur.}\sol $h=350/9.81=30.6$ m; $q=0.95\pi(30.6/2000)^2\sqrt{2\times9.81\times4.0}\times6\times10^4=17.5$ L\,min$^{-1}$ \ans $r=q\times60/(S_lS_m)=17.5\times60/216=4.9$ mm\,h$^{-1}$ \ans The rate exceeds $K_s$: ponding after $7.5$ h; sets longer than that run off, so use a smaller nozzle or a wider spacing. \ans

}{}
\IfFileExists{extra/tenbx13.tex}{\section{Worked solutions on ten crops (cool climate)}
The ten crops of the library differ in their demand, their roots and their sensitivity. The same type of question is answered for each of them at $ET_0=4.0$ mm\,d$^{-1}$; the data are typical values of the crops (Doorenbos--Kassam and FAO 56).

\exbox{Crop example 13b: Maize}{Maize: a solid-set sprinkler system with $r=7.6$ mm\,h$^{-1}$ irrigates every 3 days at $E_a=0.75$. Find the gross depth, the set time and the number of sets per 20-h day.}\sol Net depth $3\times4.80=14.4$ mm; gross $14.4/0.75=19.2$ mm \ans $t_{\rm set}=19.2/7.6=2.5$ h \ans Sets per day $\lfloor20/2.5\rfloor=7$ \ans

\exbox{Crop example 13b: Wheat}{Wheat: a solid-set sprinkler system with $r=7.6$ mm\,h$^{-1}$ irrigates every 3 days at $E_a=0.75$. Find the gross depth, the set time and the number of sets per 20-h day.}\sol Net depth $3\times4.60=13.8$ mm; gross $13.8/0.75=18.4$ mm \ans $t_{\rm set}=18.4/7.6=2.4$ h \ans Sets per day $\lfloor20/2.4\rfloor=8$ \ans

\exbox{Crop example 13b: Cotton}{Cotton: a solid-set sprinkler system with $r=7.6$ mm\,h$^{-1}$ irrigates every 3 days at $E_a=0.75$. Find the gross depth, the set time and the number of sets per 20-h day.}\sol Net depth $3\times4.60=13.8$ mm; gross $13.8/0.75=18.4$ mm \ans $t_{\rm set}=18.4/7.6=2.4$ h \ans Sets per day $\lfloor20/2.4\rfloor=8$ \ans

\exbox{Crop example 13b: Tomato}{Tomato: a solid-set sprinkler system with $r=7.6$ mm\,h$^{-1}$ irrigates every 3 days at $E_a=0.75$. Find the gross depth, the set time and the number of sets per 20-h day.}\sol Net depth $3\times4.60=13.8$ mm; gross $13.8/0.75=18.4$ mm \ans $t_{\rm set}=18.4/7.6=2.4$ h \ans Sets per day $\lfloor20/2.4\rfloor=8$ \ans

\exbox{Crop example 13b: Potato}{Potato: a solid-set sprinkler system with $r=7.6$ mm\,h$^{-1}$ irrigates every 3 days at $E_a=0.75$. Find the gross depth, the set time and the number of sets per 20-h day.}\sol Net depth $3\times4.60=13.8$ mm; gross $13.8/0.75=18.4$ mm \ans $t_{\rm set}=18.4/7.6=2.4$ h \ans Sets per day $\lfloor20/2.4\rfloor=8$ \ans

\exbox{Crop example 13b: Soybean}{Soybean: a solid-set sprinkler system with $r=7.6$ mm\,h$^{-1}$ irrigates every 3 days at $E_a=0.75$. Find the gross depth, the set time and the number of sets per 20-h day.}\sol Net depth $3\times4.60=13.8$ mm; gross $13.8/0.75=18.4$ mm \ans $t_{\rm set}=18.4/7.6=2.4$ h \ans Sets per day $\lfloor20/2.4\rfloor=8$ \ans

\exbox{Crop example 13b: Sugar beet}{Sugar beet: a solid-set sprinkler system with $r=7.6$ mm\,h$^{-1}$ irrigates every 3 days at $E_a=0.75$. Find the gross depth, the set time and the number of sets per 20-h day.}\sol Net depth $3\times4.80=14.4$ mm; gross $14.4/0.75=19.2$ mm \ans $t_{\rm set}=19.2/7.6=2.5$ h \ans Sets per day $\lfloor20/2.5\rfloor=7$ \ans

\exbox{Crop example 13b: Sunflower}{Sunflower: a solid-set sprinkler system with $r=7.6$ mm\,h$^{-1}$ irrigates every 3 days at $E_a=0.75$. Find the gross depth, the set time and the number of sets per 20-h day.}\sol Net depth $3\times4.40=13.2$ mm; gross $13.2/0.75=17.6$ mm \ans $t_{\rm set}=17.6/7.6=2.3$ h \ans Sets per day $\lfloor20/2.3\rfloor=8$ \ans

\exbox{Crop example 13b: Onion}{Onion: a solid-set sprinkler system with $r=7.6$ mm\,h$^{-1}$ irrigates every 3 days at $E_a=0.75$. Find the gross depth, the set time and the number of sets per 20-h day.}\sol Net depth $3\times4.20=12.6$ mm; gross $12.6/0.75=16.8$ mm \ans $t_{\rm set}=16.8/7.6=2.2$ h \ans Sets per day $\lfloor20/2.2\rfloor=9$ \ans

\exbox{Crop example 13b: Alfalfa}{Alfalfa: a solid-set sprinkler system with $r=7.6$ mm\,h$^{-1}$ irrigates every 3 days at $E_a=0.75$. Find the gross depth, the set time and the number of sets per 20-h day.}\sol Net depth $3\times4.60=13.8$ mm; gross $13.8/0.75=18.4$ mm \ans $t_{\rm set}=18.4/7.6=2.4$ h \ans Sets per day $\lfloor20/2.4\rfloor=8$ \ans

}{}
\IfFileExists{extra/ex13.tex}{\backhead{EXERCISES}{khred}
\begin{exlist}
\item Discharge of a 5-mm nozzle at 300 kPa.\exans{1.64 m$^3$/h}
\item Application rate for $q=1.8$ m$^3$/h, $12\times12$.\exans{12.5 mm/h}
\item Set time for 40 mm at 8 mm/h.\exans{5 h}
\item Spacing at 60 per cent of a 36-m diameter.\exans{21.6 m}
\item $CU$ for depths 10, 12, 8, 10.\exans{90.0 per cent}
\item Gross depth for $ET=6$, 4 d, $E_a=0.75$.\exans{32 mm}
\item Name two continuous-move systems.\exans{pivot, linear move}
\item Why is wind a problem?\exans{drift distorts pattern}
\item What is the 20 per cent rule?\exans{lateral pressure variation under 20 per cent}
\item Number of sprinklers for 8 ha at $12\times12$.\exans{556}
\end{exlist}
}{}
\IfFileExists{extra/dyk13.tex}{\dykband{DO YOU KNOW ?}
\begin{exlist}
\item Name three sprinkler system types.
\item What sets the wetted radius?
\item What is Christiansen's coefficient?
\item Why must $I$ be below the intake rate?
\item What is the 20 per cent rule?
\item How does wind change the spacing?
\item Why does a pivot's rate rise to the end?
\item How do you find the set time?
\end{exlist}
}{}
\IfFileExists{extra/obj13.tex}{\dykband{OBJECTIVE TYPE QUESTIONS}
\begin{exlist}
\item $I=$\opt{$q/(S_lS_m)$}{$qS$}{$q+S$}{none}
\item CU target\opt{85 per cent}{10}{20}{none}
\item Pivot is\opt{continuous-move}{basin}{furrow}{none}
\item Wind\opt{distorts}{helps}{none}{rain}
\item Orifice law uses\opt{$\sqrt{2gH}$}{$gH$}{$H$}{none}
\item Spacing is a fraction of\opt{diameter}{depth}{time}{none}
\item Solid set has\opt{fixed laterals}{no pipes}{no pumps}{none}
\item $I$ below\opt{intake}{rain}{ET}{none}
\item Set time = depth/\opt{rate}{time}{area}{none}
\item Sprinkle loses\opt{drift}{nothing}{none}{gold}
\end{exlist}
\objkey{1.~(a)\quad 2.~(a)\quad 3.~(a)\quad 4.~(a)\quad 5.~(a)\quad 6.~(a)\quad 7.~(a)\quad 8.~(a)\quad 9.~(a)\quad 10.~(a)}
}{}
\swaflvideo{13}
